% This is part of Le Frido and Giulietta
% Copyright (c) 2008-2026
%   Laurent Claessens, Carlotta Donadello
% See the file fdl-1.3.txt for copying conditions.

%+++++++++++++++++++++++++++++++++++++++++++++++++++++++++++++++++++++++++++++++++++++++++++++++++++++++++++++++++++++++++++
\section{Connexité}
%+++++++++++++++++++++++++++++++++++++++++++++++++++++++++++++++++++++++++++++++++++++++++++++++++++++++++++++++++++++++++++

L'idée de la connexité, c'est de s'assurer qu'un ensemble est «d'un seul tenant».

\begin{propositionDef}[Espace topologique connexe\cite{BIBooBQKKooUSSngj}]	\label{PROPooNLXDooVTZLan}
	Soit un espace topologique \( X\). Les propositions suivantes sont équivalentes :
	\begin{enumerate}
		\item
		      \( X\) n'est pas la réunion de deux ouverts non vides disjoints.
		\item
		      \( X\) n'est pas la réunion de deux fermés non vides disjoints.
		\item
		      Les seuls ouverts-fermés de \( X\) sont \( X\) et \( \emptyset\).
		\item
		      Toute application continue \( X\to\{ 1,2 \}\) (munie de la topologie discrète) est constante.
	\end{enumerate}
	Quand un espace topologique vérifie ces conditions, nous disons que \( X\) est \defe{connexe}{espace connexe}.
	%TODOooSRAFooVCjJiR. Prouver ça.
\end{propositionDef}

\begin{definition}[Partie connexe\cite{BIBooBQKKooUSSngj}]	\label{DEFooQGESooKSVtmt}
	Une partie d'un espace topologique est \defe{connexe}{partie connexe} si elle est connexe pour la topologie induite.
\end{definition}

\begin{definition}[Partie o-connexe\cite{MonCerveau}]	\label{DEFooJYCDooGNEDoC}
	Une partie \( Q\) d'un espace topologique \( X\) est \defe{o-connexe}{o-connexe} si pour tout ouverts \( A,B\) tels que
	\begin{enumerate}
		\item
		      \( Q\subset A\cup B\)
		\item
		      \( A\cap Q\neq\emptyset\) et \( B\cap Q\neq \emptyset\),
	\end{enumerate}
	nous avons \( A\cap B\neq \emptyset\).
\end{definition}


\begin{proposition}[\cite{MonCerveau}]	\label{PROPooQHJXooMNrxOJ}
	Tout connexe est o-connexe.
\end{proposition}

\begin{proof}
	Soient un espace topologique \( (X,\tau_X)\) ainsi qu'une partie \( Q\subset X\) munie de la topologie induite \( \tau_Q\). Nous supposons que \( Q\) n'est pas o-connexe, et nous montrons que \( Q\) n'est pas connexe.

	Si \( Q\) n'est pas o-connexe, il existe \( A,B\in \tau_X\) tels que \(  Q\subset A\cup B\), \( A\cap B=\emptyset\), \( A\cap Q\neq\emptyset\neq B\cap Q\).

	Les parties \( A\cap Q\) et \( B\cap Q\) sont des ouverts de \( (Q,\tau_Q)\). Elles sont disjointes, non vides et recouvrent \( Q\). Donc \( Q\) n'est pas connexe.
\end{proof}

\begin{proposition}[\cite{BIBooCJEYooZiHPbC}]	\label{PROPooAHWNooXhrqnd}
	Dans un espace complètement normal\footnote{Définition \ref{DEFooTLVHooYWvDIu}.}, tout o-connexe est connexe.
\end{proposition}

\begin{proof}
	Soient un espace topologique complètement normal \( X\), et une partie \( Q\in X\).

	\begin{subproof}
		\spitem[Contraposée]
		%-----------------------------------------------------------
		Par contraposée nous supposons que \( Q\) n'est pas connexe : il existe \( G_0,G_1\in \tau_Q\) tels que
		\begin{enumerate}
			\item
			      \( Q=G_0\cup G_1\)
			\item
			      \( G_0\neq\emptyset\neq G_1\)
			\item
			      \( G_0\cap G_1\neq \emptyset\)
		\end{enumerate}
		Par définition de la topologie induite \( \tau_Q\), il existe des ouverts \( U_0,U_1\in \tau_X\) tels que \( G_0=U_0\cap Q\) et \( G_1=U_1\cap Q\).
		\spitem[\( U_0\cap G_1=\emptyset\)]
		%-----------------------------------------------------------
		Nous avons
		\begin{equation}
			U_0\cap G_1=U_0\cap U_1\cap Q\subset U_0\cap Q=G_0.
		\end{equation}
		Donc un élément de \( U_0\cap G_1\) serait dans \( G_0\cap G_1\), ce qui n'est pas possible.

		\spitem[Fermeture]
		%-----------------------------------------------------------
		Vu que \( U_0\) est ouvert dans \( X\) et que \( U_0\cap G_1=\emptyset\), le lemme \ref{LEMooLKHEooRqrqkA} dit que \( U_0\cap \Adh_X(G_1)=\emptyset\). Et comme \( G_0\subset U_0\),
		\begin{equation}
			G_0\cap\Adh_X(G_1)=\emptyset.
		\end{equation}
		De même
		\begin{equation}
			G_1\cap\Adh_X(G_0)=\emptyset.
		\end{equation}

		\spitem[Complètement normal]
		%-----------------------------------------------------------
		Vu que \( X\) est complètement normal, il existe des ouverts disjoints \( V_0,V_1\in \tau_X\) tels que \( G_0\subset V_0\) et \( G_1\subset V_1\).

		\spitem[Conclusion]
		%-----------------------------------------------------------
		Nous concluons que \( Q\) n'est pas o-connexe parce que \( \{ V_0,V_1 \}\) est un recouvrement disjoint de \( Q\) par des ouverts de \( X\).
	\end{subproof}
\end{proof}

\begin{probleme}
	Il y a peut-être une réciproque. J'ai entendu à quelques endroits sur internet que si tous les o-connexes sont connexes, alors l'espace est complètement normal. Voici une tentative de démonstration, pas complète. Si vous pouvez donner une démonstration complète ou un contre-exemple (un espace non complètement normal dans lequel tous les o-connexes sont connexes), écrivez-moi.
\end{probleme}

\begin{pseudotheorem}[\cite{BIBooCJEYooZiHPbC,BIBooPGBBooDodAGg}]	\label{THOooMULNooNYeHAj}
	Un espace topologique dans lequel tous les o-connexes sont connexes est complètement normal.
\end{pseudotheorem}

% jour on a une preuve, il faudra supprimer ooVQTCooVAMjhs parce que la ligne contient une référence vers le futur que j'ai accepté parce que de toutes façons c'est pas une vraie preuve. ooFFZUooSPBhvW

\begin{proof}
	Soient des parties \( A,B\subset X\)  telles que \( A\cap\bar B=B\cap\bar A=\emptyset\). Nous notons \( Q=A\cup B\).
	\begin{subproof}
		\spitem[\( B\) est ouvert dans \(Q= A\cup B\)]
		%-----------------------------------------------------------
		Nous savons que \( X\setminus \bar B\) est ouvert dans \( X\). Vu que \( B\cap\bar A=\emptyset\) nous avons \( B\subset Q\cap(X\setminus\bar A)\), et donc \( B\subset X\setminus \bar A\) et donc \( B=Q\cap(X\setminus \bar A)\). La partie \( B\) est donc bien l'intersection entre \( Q\) et un ouvert de \( X\). Donc \( B\) est ouvert dans \( Q\).

		Donc \( A\cup B\) est l'union de deux ouverts disjoints (pour sa topologie induite), et n'est donc pas connexe.

		\spitem[Utilisation de l'hypothèse]
		%-----------------------------------------------------------
		Vu que \( Q\) n'est pas connexe, il n'est pas non plus o-connexe. Il existe des ouverts \( U\) et \( V\) tels que
		\begin{enumerate}
			\item
			      \( Q\subset U\cup V\)
			\item
			      \( U\cap Q\neq\emptyset\neq V\cap Q\)
			\item
			      \( U\cap V=\emptyset\).
		\end{enumerate}

		Ce serait bien que \( A\subset U\) et que \( B\subset V\). Mais rien n'est moins sûr. Nous considérons donc quelques cas possibles.

		\spitem[Si \( A\) est connexe]
		%-----------------------------------------------------------
		Dans ce cas, \( \{ U,V \}\) est un recouvrement de \( A\) par des ouverts disjoints. Donc \( A\subset U\) et \( A\cap V=\emptyset\) (ou le contraire).

		\spitem[Si \( A\) a un nombre fini de composantes connexes]
		%-----------------------------------------------------------
		Soient \( \{ C_i \}_{1=1,\ldots,n}\) les composantes connexes\footnote{Définition \ref{DEFooFHXNooJGUPPn}.} de \( A\). Pour chaque \( i\) nous avons \( C_1\cap\bar B=B\cap\bar C_i=\emptyset\) et donc (en reprenant toute la démonstration jusqu'ici avec \( C_i\) au lieu de \( A\)), nous avons des ouverts \( U_i\) et \( V_i\) tels que \( C_i\subset U_i\), \( B\subset V_i\) et \( U_i\cap V_i= \emptyset\). % ooVQTCooVAMjhs ne pas enlever ceci parce que ça corrige une référence vers le futur. Voir ooFFZUooSPBhvW.

		Nous posons alors \( U=\bigcup_{i=1}^n C_i\) et \( V=\bigcap_{i=1}^nV_i\).
		\begin{subproof}
			\spitem[\( V\cap U=\emptyset\)]
			%-----------------------------------------------------------
			En effet si \( x\in V\cap U\) nous avons \( V\in C_i\) pour un certain \( i\). Pour ce même \( i\) nous avons \( x\in V_i\) (c'est pour ça qu'on a dit que \( V\) serait l'intersection). Or \( V_i\cap U_i=\emptyset\); un tel \( x\) ne peut pas exister.

			\spitem[\( V\neq \emptyset\)]
			%-----------------------------------------------------------
			Parce que chacun des \( V_i\) contient \( B\).

			\spitem[\( V\) est ouvert]
			%-----------------------------------------------------------
			Parce que c'est une intersection finie d'ouverts\footnote{C'est la propriété \ref{DefTopologieGene}\ref{ITEMooGJAPooOWIOzv}.}.
		\end{subproof}
		\spitem[Le problème]
		%-----------------------------------------------------------
		Si \( A\) et \( B\) ont tous les deux un nombre infini de composantes connexes, on ne peut pas déduire que \( U\) et \( V\) sont ouverts. Donc le démonstration coince dans ce cas.
	\end{subproof}

	Bref. Preuve pas terminée. Écrivez-moi si vous êtes capable de terminer.

\end{proof}

\begin{definition}  \label{DefIRKNooJJlmiD}
	Lorsque \( X\) est un espace topologique, nous disons qu'un sous-ensemble \( A\) est \defe{non connexe}{connexité!définition} quand on peut trouver des ouverts \( O_1\) et \( O_2\) disjoints tels que
	\begin{equation}    \label{EqDefnnCon}
		A=(A\cap O_1)\cup (A\cap O_2),
	\end{equation}
	et tels que \( A\cap O_1\neq\emptyset\), et \( A\cap O_2\neq\emptyset\). Si un sous-ensemble n'est pas noo-connexe, alors on dit qu'il est \defe{connexe}{ensemble connexe}.
	%TODOooSRAFooVCjJiR. Cette définition doit être virée.
\end{definition}
Une autre façon d'exprimer la condition \eqref{EqDefnnCon} est de dire que \( A\) n'est pas connexe quand il est contenu dans la réunion de deux ouverts disjoints qui intersectent tous les deux \( A\).


\begin{proposition}[\cite{BIBooCYZRooAHkNGu}]		\label{PROPooBCFXooRlMvch}
	Soit un espace topologique \( X\). Soient \( (C_i)_{i\in I}\) des connexes de \( X\) tels que \( \bigcap_{i\in I}C_i\neq \emptyset\). Alors \( \bigcup_{i\in I}C_i\) est connexe.
\end{proposition}

\begin{proof}
	Nous notons \( C=\bigcup_{i\in I}C_i\). Soient des ouverts disjoints \( A,B\in \tau_C\) recouvrant \( C\). Nous devons démontrer que soit \( A\cap C\) soit \( B\cap C\) est vide. Par définition de la topologie induite \( \tau_C\), il existe des ouverts \( A',B'\in\tau_X\) tels que \( A=C\cap A'\) et \( B=C\cap B'\).

	Nous notons \( A_i=C_i\cap A\) et \( B_i=C_i\cap B\). Les parties \( A_i\) et \( B_i\) recouvrent \( C_i\) et nous avons \( A_i\cup B_i=C_i\). De plus
	\begin{equation}
		A_i=C_i\cap^C\cap A'=C_i\cap A',
	\end{equation}
	donc \( A_i\) est ouvert dans \( C_i\). Vu que \( C_i\) est ouvert, soit \( A_i\) soit \( B_i\) est vide.

	Par hypothèse, \( \bigcap_{i\in I}C_i\) est non vide. Soit \( x\) un élément de cette intersection. Soit \( x\in A\) et alors \( x\in A_i\) pour tout \( i\); soit \( x\in B\) et alors \( x\in B_i\) pour tout \( i\).

	Dans le premier cas, \( x\in A_i\) pour tout \( i\), de telle sorte que \( B_i=\emptyset\) pour tout \( i\). Nous avons alors
	\begin{equation}
		\emptyset=\bigcup_{i\in I}B_i=\bigcup_{i\in I}(C_i\cap B)=B\cap C.
	\end{equation}
	Dans le second cas nous obtenons de même que \( A\cap C=\emptyset\).
\end{proof}


\begin{propositionDef}[\cite{BIBooCGIFooGvxBWL}]        \label{DEFooFHXNooJGUPPn}
	Soient un espace topologique \( X\) et un point \( x\in X\).
	\begin{enumerate}
		\item		\label{ITEMooBZAQooNwuzaS}
		      La réunion de toutes les parties connexes de \( X\) contenant \( x\) est connexe.
		\item
		      Cette réunion est la plus grande (au sens de la relation d'inclusion) de toutes les parties connexes de \( X\) contenant \( x\).
	\end{enumerate}
	La réunion de toutes les parties connexes de \( X\) contenant \( x\) est nommée \defe{composante connexe}{composante connexe} de \( x\) dans \( X\).
\end{propositionDef}

\begin{proof}
	La réunion de toutes les parties connexes contenant \( x\) est connexe par la proposition \ref{PROPooBCFXooRlMvch}.

	Nous notons \( C\) la réunion de toutes les parties connexes contenant \( x\). Si \( D\) est un connexe contenant \( x\), alors \( D\subset C\) parce que \( C\) est une union de tous les connexes, y compris \( D\).
\end{proof}

\begin{proposition}[\cite{MonCerveau}] \label{PropHSjJcIr}
	Soit \( X\) un espace topologique. Les conditions suivantes sont équivalentes.
	\begin{enumerate}
		\item       \label{ITEMooXHIKooGqrgTs}
		      L'espace \( X\) est connexe.
		\item       \label{ITEMooRTNPooADKVnw}
		      Si \( X=O_1\cup O_2\) avec \( O_1\) et \( O_2\) des ouverts disjoints, alors soit \( O_1=\emptyset\) soit \( O_2=\emptyset\).
		\item       \label{ITEMooOEZYooFBNaOZ}
		      Si \( X=F_1\cup F_2\) avec \( F_1\) et \( F_2\) fermés disjoints dans \( X\), alors \( F_1=\emptyset\) ou \( F_2=\emptyset\).
		\item       \label{ITEMooNIPZooIDPmEf}
		      Si \( A\subset X\) avec \( A\) ouvert et fermé en même temps, alors \( A=\emptyset\) ou \( A=X\).
	\end{enumerate}
\end{proposition}

\begin{proof}
	En quatre parties.
	\begin{subproof}
		\spitem[\ref{ITEMooXHIKooGqrgTs} implique \ref{ITEMooRTNPooADKVnw}]
		Par rapport à la définition \ref{DefIRKNooJJlmiD}, nous prenons la partie \( X\) de l'espace \( X\). Supposons que \( O_1\) et \( O_2\) sont tout deux non vides. Dans ce cas nous avons
		\begin{equation}
			X= O_1\cup O_2 = (O_1\cap X)\cup (O_2\cap X),
		\end{equation}
		ce qui prouverait que \( X\) est non connexe. Contradiction. Un des \( O_i\) est vide.
		\spitem[\ref{ITEMooRTNPooADKVnw} implique \ref{ITEMooOEZYooFBNaOZ}]
		Soit une union disjointe de fermés \( X=F_1\cup F_2\). Puisque l'union est disjointe, nous avons \( F_1=X\setminus F_2\) et \( F_2=X\setminus F_1\), ce qui fait que \( F_1\) et \( F_2\) sont également ouverts. Nous en déduisons que \( X=F_1\cup F_2\) est une union disjointe d'ouverts. L'hypothèse indique que \( F_1=\emptyset\) ou \( F_2=\emptyset\).
		\spitem[\ref{ITEMooOEZYooFBNaOZ} implique \ref{ITEMooNIPZooIDPmEf}]
		Soit \( A\) une partie ouverte et fermée de \( X\). Nous supposons que \( A\) est ouvert et fermé, donc \( X\setminus A\) est également ouvert et fermé : c'est la définition \ref{DEFFermeooNSAAooHxZbAo} d'un fermé. Nous avons évidemment l'union \( X=A\cup(X\setminus A)\) qui est une union disjointe de fermés. Par hypothèse nous avons soit \( A=\emptyset\) soit \( X\setminus A=\emptyset\).
		\spitem[\ref{ITEMooNIPZooIDPmEf} implique \ref{ITEMooXHIKooGqrgTs}]
		Supposons que \( X\) ne soit pas connexe. Il existe donc des ouverts disjoints \( O_1\) et \( O_2\) tels que \( X=O_1\cup O_2\). Étant donné que \( O_1=X\setminus O_2\), la partie \( O_1\) est fermée comme complémentaire d'ouvert. Donc \( O_1\) est fermé et ouvert (et \( O_2\) aussi d'ailleurs). Par hypothèse nous concluons que \( O_1\) est soit \( X\) soit \( \emptyset\).
	\end{subproof}
\end{proof}

Nous verrons plus tard (proposition~\ref{PropConnexiteViaFonction}) une autre caractérisation de la connexité basée sur la continuité des fonctions \( X\to \eZ\).

\begin{proposition}[\cite{MonCerveau}]     \label{PROPooSCKNooRbewdv}
	Soient un espace topologique \( X\) ainsi que \( S\subset X\). Si \( U\subset X\) est connexe\footnote{Définition \ref{DefIRKNooJJlmiD}.} et si \( U\subset S\subset \bar U\), alors \( S\) est connexe.
\end{proposition}

\begin{proof}
	Supposons que \( S\) ne soit pas connexe. Il existe des ouverts disjoints \( A\) et \( B\) tels que \( S\subset A\cup B\) et \( S\cap A\neq\emptyset\), \( S\cap B\neq\emptyset\). Nous prouvons que ces ouverts fonctionnent aussi pour prouver que \( U\) est non connexe (donc on aura une contradiction).

	D'abord \( A\) et \( B\) recouvrent \( U\) parce que \( U\subset S\subset A\cup B\). Ensuite prouvons que \( U\cap A\neq \emptyset\). Soit \( x\in S\cap A\). Vu que \( x\in S\subset\bar U\), tout voisinage de \( x\) intersecte \( U\) (c'est la définition \ref{DEFooSVWMooLpAVZR} de l'adhérence). De plus \( x\in A\) et \( A\) est ouvert. Soit donc un voisinage \( V\) de \( x\) contenu dans \( A\). Nous avons \( V\subset A\) et \( V\cap U\neq \emptyset\). Donc \( A\cap U\neq \emptyset\).

	Le même raisonnement tient pour \( B\).
\end{proof}

\begin{proposition} \label{PropIWIDzzH}
	Stabilité de la connexité par union.
	\begin{enumerate}
		\item       \label{ITEMooLVSSooTGstBz}
		      Une union quelconque de connexes ayant une intersection non vide est connexe.
		\item
		      Pour tout \( n \in \eN, n > 0 \), si \( A_1,\ldots, A_n\) sont des connexes de \( X\) avec \( A_i\cap A_{i+1}\neq \emptyset\), alors l'union \( \bigcup_{i=1}^nA_i\) est connexe.
	\end{enumerate}
\end{proposition}

\begin{proof}
	Point par point.
	\begin{enumerate}
		\item
		      Soient \( \{ C_i \}_{i\in I}\) un ensemble de connexes et un point \( p\) dans l'intersection : \( p\in\bigcap_{i\in I}C_i\). Supposons que l'union ne soit pas connexe. Alors nous considérons \( A\) et \( B\), deux ouverts disjoints recouvrant tous les \( C_i\) et ayant chacun une intersection non vide avec l'union.

		      Supposons pour fixer les idées que \( p\in A\) et prenons \( x\in B\cap\bigcup_{i\in I}C_i\). Il existe un \( j\in I\) tel que \( x\in C_j\). Avec tout cela nous avons
		      \begin{enumerate}
			      \item
			            \( C_j\subset A\cup B\) parce que \(A \cup B\) recouvre tous les \( C_i \),
			      \item
			            \( C_j\cap A\neq \emptyset\) parce que \( p\) est dans l'intersection,
			      \item
			            \( C_j\cap B\neq\emptyset\) parce que \( x\) est dans cette intersection.
		      \end{enumerate}
		      Cela contredit le fait que \( C_j\) soit connexe.

		\item

		      Pour la seconde partie nous procédons de proche en proche\footnote{Parce qu'on a la flemme de rédiger correctement une récurrence.}. D'abord \( A_1\cup A_2\) est connexe par la première partie, ensuite \( (A_1\cup A_2)\cup A_3\) est connexe parce que les connexes \( A_1\cup A_2\) et \( A_3\) ont un point d'intersection par hypothèse, et ainsi de suite.
	\end{enumerate}
\end{proof}


%+++++++++++++++++++++++++++++++++++++++++++++++++++++++++++++++++++++++++++++++++++++++++++++++++++++++++++++++++++++++++++
\section{Compacité}
%+++++++++++++++++++++++++++++++++++++++++++++++++++++++++++++++++++++++++++++++++++++++++++++++++++++++++++++++++++++++++++

La compacité est le thème~\ref{THEMEooQQBHooLcqoKB}.

%---------------------------------------------------------------------------------------------------------------------------
\subsection{Définition et notions connexes}
%---------------------------------------------------------------------------------------------------------------------------

Soit \( E\), un sous-ensemble de \( \eR\). Nous pouvons considérer les ouverts suivants :
\begin{equation}
	\mO_x=B(x,1)
\end{equation}
pour chaque \( x\in E\). Évidemment,
\begin{equation}
	E\subseteq \bigcup_{x\in E}\mO_x.
\end{equation}
Cette union contient en général de nombreuses redondances. Si par exemple \( E=[-10,10]\), l'élément \( 3\in E\) est contenu dans \( \mO_{3.5}\), \( \mO_{2.7}\) et bien d'autres. Pire : même si on enlève par exemple \( \mO_2\) de la liste des ouverts, l'union de ce qui reste continue à être tout \( E\). La question est : \emph{est-ce qu'on peut en enlever suffisamment pour qu'il n'en reste qu'un nombre fini ?}

\begin{definition}
	Soit \( E\), un sous-ensemble de \( \eR\). Une collection d'ouverts \( \mO_i\) est un \defe{recouvrement}{recouvrement} de \( E\) si \( E\subseteq \bigcup_{i}\mO_i\).
\end{definition}

\begin{definition} \label{DefJJVsEqs}
	Une partie \( A\) d'un espace topologique est \defe{compacte}{compact} si elle vérifie la propriété de Borel-Lebesgue : pour tout recouvrement de \( A\) par des ouverts (c'est-à-dire une collection d'ouverts dont la réunion contient \( A\)) on peut extraire un recouvrement fini.
\end{definition}


\begin{remark}
	Certaines sources (dont \wikipedia{fr}{Compacité_(mathématiques)}{Wikipédia}) disent que pour être compact il faut aussi être séparé\footnote{Définition~\ref{DefYFmfjjm}.}. Pour ces sources, un espace qui ne vérifie que la propriété de Borel-Lebesgue est alors dit \defe{quasi-compact}{quasi-compact}\index{compact!quasi}.
\end{remark}

\begin{normaltext}
	La définition \ref{DefJJVsEqs} en cache deux. En effet, si la partie \( A\) est l'espace topologique lui-même, cela définit un espace topologique compact. Un espace topologique est compact \emph{en soi} lorsque de tout recouvrement par des ouverts, nous pouvons extraire un sous-recouvrement fini. Dans ce cas, si \( X\) est l'espace et si \( \{ A_i \}_{i\in I}\) est le recouvrement, nous avons \( X=\bigcup_{i\in I}A_i\) et non une simple inclusion \( X\subset \bigcup_{i\in I}A_i\).
\end{normaltext}

\begin{lemma}       \label{LEMooVYTRooKTIYdn}
	Si \( K\) est une partie compacte de l'espace topologique \( X\), alors \( K\) est un espace topologique compact pour la topologie induite\footnote{Définition \ref{DefVLrgWDB}.} de \( X\).
\end{lemma}

\begin{proof}
	Nous notons \( \tau\) la topologie de \( X\) et \( \tau_K\) la topologie induite de \( X\) vers \( K\), c'est-à-dire
	\begin{equation}
		\tau_K=\{ \mO\cap K\tq \mO\in\tau \}.
	\end{equation}
	Soient des ouverts \( A_i\in \tau_K\) (\( i\in I\) où \( I\) est un ensemble quelconque) tels que \( \bigcup_iA_i=K\). Pour chaque \( i\in I\), il existe un \( \mO_i\in \tau\) tel que \( A_i=K\cap\mO_i\). Nous avons
	\begin{equation}
		K=\bigcup_{i\in I}(K\cap\mO_i)\subset\bigcup_{i\in I}\mO_i.
	\end{equation}
	Donc les \( \mO_i\) forment un recouvrement de \( K\) par des ouverts de \( X\). Puisque \( K\) est une partie compacte de \( X\), il existe un sous-ensemble fini \( J\) de \( I\) tel que
	\begin{equation}
		K\subset\bigcup_{j\in J}\mO_j.
	\end{equation}
	Nous avons donc aussi
	\begin{equation}
		K\subset\bigcup_{j\in J}K\cap\mO_j=\bigcup_{j\in J}A_j.
	\end{equation}
	Nous avons prouvé que \( \{ A_j \}_{j\in J}\) est un recouvrement fini de \( K\) par des ouverts de \( K\). Donc \( K\) est un espace topologique compact.
\end{proof}

\begin{proposition}[\cite{BIBooQKARooMHqitK}]       \label{PROPooOXKSooEDzCRZ}
	Soit \( K\) compact dans \( \eC\). Si \( \mO\) est une composante connexe\footnote{Définition \ref{DEFooFHXNooJGUPPn}.} de \( \eC\setminus K\), alors \( \partial \mO\subset K\).
\end{proposition}

\begin{proof}
	Supposons que \( \partial\mO\) n'est pas inclus dans \( K\), et considérons \( w\in\partial\mO\setminus K\). Nous posons \( A=\mO\cup\{ w \}\). Vu que \( w\in\partial\mO\), nous avons \( w\in\bar\mO\). Donc
	\begin{equation}
		\mO\subset A\subset\bar\mO.
	\end{equation}
	La proposition \ref{PROPooSCKNooRbewdv} fait que \( A\) est connexe parce que \( \mO\) est connexe. Donc \( A\) est un connexe de \( \eC\setminus K\) contenant strictement \( \mO\). Impossible parce que \( \mO\) est une composante connexe de \( \eC\setminus K\). Contradiction.

	Nous en déduisons que \( \partial\mO\subset K\).
\end{proof}



%---------------------------------------------------------------------------------------------------------------------------
\subsection{Autres compacité}
%---------------------------------------------------------------------------------------------------------------------------

\begin{definition}[Séquentiellement compact]        \label{DEFooTVDOooZbwOFK}
	Nous disons qu'un espace topologique est \defe{séquentiellement compact}{compact!séquentiellement} si toute suite admet une sous-suite convergente.
\end{definition}

\begin{definition}[\cite{BIBooYNXFooXgzmvu}]	\label{DEFooSOLWooWlSiUn}
	Un espace topologique est \defe{\( \sigma\)-compact}{\( \sigma\)-compact} si il est union dénombrable de sous-espaces compacts.

	On dit aussi \defe{dénombrable à l'infini}{dénombrable à l'infini}.
\end{definition}

\begin{definition}
	Une famille \( \mA\) de parties de \( X\) a la \defe{propriété d'intersection finie non vide}{propriété d'intersection non vide} si tout sous-ensemble fini de \( \mA\) a une intersection non vide.
\end{definition}

\begin{proposition}\label{PropXKUMiCj}
	Soient \( X\) un espace topologique et \( K\subset X\). Les propriétés suivantes sont équivalentes :
	\begin{enumerate}
		\item\label{ItemXYmGHFai}
		      \( K\) est compact.
		\item\label{ItemXYmGHFaii}
		      Si \( \{ F_i \}_{i\in I}\) est une famille de fermés telle que \( \bigcap_{i\in I}F_i \cap K =\emptyset\), alors il existe une partie finie non vide \( A\) de \( I\) tel que \( \bigcap_{i\in A}F_i \cap K =\emptyset\).
		\item\label{ItemXYmGHFaiii}
		      Si \( \{ F_i \}_{i\in I}\) est une famille de fermés telle que pour tout choix de \( A\) fini dans \( I\) nous ayons \( \bigcap_{i\in A}F_i \cap K \neq\emptyset\), alors l'intersection complète est non vide : \( \bigcap_{i\in I}F_i \cap K\neq\emptyset\).
		\item\label{ItemXYmGHFaiv}
		      Toute famille de fermés de \( X \), à laquelle \( K \) est joint, et qui a la propriété d'intersection finie non vide, a une intersection non vide.
	\end{enumerate}
\end{proposition}

\begin{proof}
	Les propriétés~\ref{ItemXYmGHFaiii} et~\ref{ItemXYmGHFaii} sont équivalentes par contraposition. De plus le point~\ref{ItemXYmGHFaiv} est une simple\footnote{Enfin, simple\dots{} il faut remarquer que dans la formulation de~\ref{ItemXYmGHFaiv}, les intersections peuvent ne pas faire intervenir \( K \), mais, au final, on s'en moque.} reformulation en français de la propriété~\ref{ItemXYmGHFaiii}.

	Prouvons~\ref{ItemXYmGHFai} \( \Rightarrow\)~\ref{ItemXYmGHFaii}. Soit \( \{ F_i \}_{i\in I}\) une famille de fermés tels que \( K\bigcap_{i\in I}F_i=\emptyset\). Les complémentaires \( \mO_i\) de \( F_i\) dans \( X\) recouvrent \( K\) et donc on peut en extraire un sous-recouvrement fini :
	\begin{equation}
		K\subset\bigcup_{i\in A}\mO_i
	\end{equation}
	pour un certain sous-ensemble fini \( A\) de \( I\). Pour ce même choix \( A\), nous avons alors aussi
	\begin{equation}
		\bigcap_{i\in A}F_i \cap K =\emptyset.
	\end{equation}

	L'implication~\ref{ItemXYmGHFaii} \( \Rightarrow\)~\ref{ItemXYmGHFai} est la même histoire de passage aux complémentaires.
\end{proof}

Le théorème \ref{ThoCQAcZxX} est en général celui qu'on nomme «théorème des fermés emboîtés», mais le corolaire suivant en mériterait également le nom.
\begin{corollary}[\cite{MonCerveau}]       \label{CORooQABLooMPSUBf}
	Soient un espace topologique compact \( X\) et une suite \( (F_i)_{i\in \eN}\) de fermés emboîtés\footnote{C'est-à-dire que \( F_{i+1}\subset F_i\).} dans \( X\) telle que
	\begin{equation}
		\bigcap_{i\in \eN}F_i=\emptyset.
	\end{equation}
	Alors il existe \( j_0\in \eN\) tel que \( F_i=\emptyset\) pour tout \( i\geq j_0\).
\end{corollary}

\begin{proof}
	La proposition \ref{PropXKUMiCj} nous dit qu'il existe une partie finie non vide \( J\) de \( \eN\) telle que \( \bigcup_{j\in J}F_j=\emptyset\). Si \( j_0=\min(J)\), alors \( F_j\subset F_{j_0}\) pour tout \( j\in J\) et nous avons
	\begin{equation}
		\emptyset=\bigcap_{j\in J}F_j=F_{j_0}.
	\end{equation}
	Dès que \( F_{j_0}=\emptyset\), tous les suivants sont également vides.
\end{proof}




%---------------------------------------------------------------------------------------------------------------------------
\subsection{Quelques propriétés}
%---------------------------------------------------------------------------------------------------------------------------

\begin{lemma}   \label{LemOWVooZKndbI}
	Une partie \( K\) d'un espace topologique est compacte si et seulement si de tout recouvrement par des ouverts d'une base de topologie nous pouvons extraire un sous-recouvrement fini.
\end{lemma}
Remarquons que la partie qui est réellement à prouver est que, si \og ça marche \fg{} pour des ouverts d'une base de topologie, alors \og ça marche\fg{} pour tous types d'ouverts.
\begin{proof}
	Soit \( K\) une partie d'un espace topologique et \( \{ \mO_i \}_{i\in I}\) un recouvrement de \( K\) par des ouverts. Chacun des \( \mO_i\) est une union d'éléments de la base de topologie par la proposition~\ref{DEFooLEHPooIlNmpi}: disons \( \mO_i = \bigcup_{j \in J_i} A_{(i,j)} \). Soit \( J = \{ j = (i, j_i) | i \in I, j_i \in J_i \} \); alors nous obtenons  \( \bigcup_{j\in J}A_j=\bigcup_{i\in I}\mO_i\).

	Par hypothèse nous pouvons extraire un ensemble fini \( J_0\subset J\) tel que \( K\subset\bigcup_{j\in J_0}A_j\). Par construction chacun des \( A_j\) est inclus dans (au moins) un des \( \mO_i\). Le choix d'un élément de \( I\) pour chacun des éléments de \( J_0\) donne une partie finie \( I_0\) de \( I\) telle que \( K\subset\bigcup_{j\in J_0}A_j\subset\bigcup_{i\in I_0}\mO_i\).
\end{proof}


\begin{example}[Un compact non fermé]
	En général, un compact n'est pas toujours fermé. Si nous prenons par exemple un ensemble \( X\) de plus de deux points muni de la topologie grossière \( \{ \emptyset,X \}\). Toutes les parties de cet espace sont compactes, mais les seuls fermés sont \( \{ \emptyset,X \}\). Toutes les autres parties sont alors compactes et non fermées.
\end{example}


\begin{lemma}[Compacts et fermés\cite{SNSposN}]   \label{LemnAeACf}
	À propos de parties fermées dans un compact.
	\begin{enumerate}
		\item       \label{ITEMooNKAKooQoNddr}
		      Une partie fermée d'un compact est compacte.
		\item       \label{ITEMooAZWVooLyPDeY}
		      Tout compact d'un espace topologique séparé est fermé.
	\end{enumerate}
\end{lemma}
\index{compacts et fermés}

\begin{proof}
	En deux parties.
	\begin{subproof}
		\spitem[Pour \ref{ITEMooNKAKooQoNddr}]
		Soient \( F\) fermé dans un compact \( K\) et \( \{ \mO_i \}_{i\in I}\) un recouvrement de \( F\) par des ouverts. Puisque \( F\) est fermé, \( F^c\) est ouvert et \( \{ \mO_i \}_{i\in I}\cup\{ K\setminus F \}\) est un recouvrement de \( K\) par des ouverts. Si nous en extrayons un sous-recouvrement fini, c'est un recouvrement de \( F\), et en supprimant éventuellement l'ouvert \( K\setminus F\), ça reste un sous-recouvrement fini de \( F\) tout en étant extrait de \( \{ \mO_i \}_{i\in I}\).

		\spitem[Pour \ref{ITEMooAZWVooLyPDeY}]
		Soient \( X\) un espace séparé et \( K\) compact dans \( X\). Nous considérons \( y \in K^c\) et, par hypothèse de séparation, pour chaque \( x\in K\) nous considérons un voisinage ouvert \( V_x\) de \( x\) et un voisinage ouvert~\footnote{Oui, la notation du voisinage peut surprendre, mais elle est quand même pratique pour ce qu'on veut en faire.} \( W_x\) de \( y\) tels que \( V_x\cap W_x=\emptyset\). Bien entendu les \( V_x\) forment un recouvrement de \( K\) par des ouverts dont nous pouvons extraire un sous-recouvrement fini : soit \( S\) fini dans \( K\) tel que
		\begin{equation}
			K\subset\bigcup_{x\in S}V_x.
		\end{equation}
		L'ensemble \( W=\bigcap_{x\in S}W_x\) est une intersection finie d'ouverts autour de \( y\) et est donc un ouvert autour de \( y\).

		Montrons que \( W\cap K=\emptyset\). Soit \( a\in K\); par définition de \( S\), il existe \( s\in S\) tel que \( a\in V_s\). Par conséquent, \( a\) n'est pas dans \( W_s\) et donc pas non plus dans \( W\).

		L'ouvert \( W\) prouve que \( y\) est dans l'intérieur du complémentaire de \( K\), et comme \( y \) est arbitraire, nous concluons que le complémentaire de \( K\) est ouvert (théorème~\ref{ThoPartieOUvpartouv}), en d'autres termes, que \( K\) est fermé.
	\end{subproof}
\end{proof}


\begin{corollary}[\cite{BIBChatGPT,BIBooCMYGooDzAKmh}]		\label{CORooSSFFooNkNmlS}
	Dans un espace topologique, une intersection d'un compact avec un fermé est compacte.
\end{corollary}

\begin{proof}
	Soit un espace topologique \( X\). Soient \( K\) compact dans \( X\) et \( F\) fermé dans \( X\). Soit un recouvrement \( \{ U_i \}_{i\in I}\) de \( K\cap F\) par des ouverts de \( X\). Vu que \( K=(K\cap F)\cup (K\setminus F)\), la famille
	\begin{equation}
		\big\{ \{ U_i \}_{i\in I},X\setminus F \big\}
	\end{equation}
	recouvre \( K\). Par compacité de \( K\), nous pouvons en extraire un sous-recouvrement fini. Il existe une partie finie \( I_0\subset I\) tel que
	\begin{equation}
		\big\{ \{ U_i \}_{i\in I_0},X\setminus F \big\}
	\end{equation}
	recouvre \( K\). Cette famille recouvre donc aussi \( K\cap F\). Mais comme \( X\setminus F\) n'intersecte surement pas \( K\cap F\), la famille \( \{ U_i \}_{i\in I_0}\) recouvre \( K\cap F\). Nous avons donc bien extrait un sous-recouvrement fini pour \( K\cap F\).
\end{proof}


\begin{proposition}[Théorème des compacts emboîtés\cite{BIBooAHEQooTUONeJ}]	\label{PROPooFXCXooXNeXWl}
	Soit un espace topologique séparé \( X\). Soient des compacts emboîtés non vides \( (K_n)_{n\in \eN}\). Alors l'intersection est compacte et non vide.
\end{proposition}

\begin{proof}
	En deux parties.
	\begin{subproof}
		\spitem[L'intersection est non vide]
		%-----------------------------------------------------------
		Nous utilisons la proposition \ref{PropXKUMiCj}\ref{ItemXYmGHFaiii}. Vu que \( K_0\) est compact, nous considérons la famille de fermés\footnote{Tout compact d'un espace séparé est fermé, lemme \ref{LemnAeACf}\ref{ITEMooAZWVooLyPDeY}.} \( \{ K_i \}_{i\in \eN}\). Si \( I\) est fini dans \( \eN\), alors
		\begin{equation}
			K_{\max(I)+1}\subset\bigcap_{i\in I}K_i,
		\end{equation}
		de telle sorte que \( \bigcap_{i\in I}K_i\) est non vide. La proposition \ref{PropXKUMiCj}\ref{ItemXYmGHFaiii} dit alors que l'intersection \( \bigcap_{i\in \eN}K_i\) est non vide.

		\spitem[L'intersection est compacte]
		%-----------------------------------------------------------
		Nous posons \( K=\bigcap_{i\in \eN}K_i\). Soit un recouvrement \( \{ A_i \}_{i\in I}\) de \( K\) par des ouverts.

		\begin{subproof}
			\spitem[\( X\setminus K\) est ouvert]
			%-----------------------------------------------------------
			Vu que \( K_i\) est compact, il est fermé et \( X\setminus K_i\) est ouvert. Le lemme \ref{LEMooHRKAooRskzQL}\ref{ITEMooYYWAooLnADjR} dit que
			\begin{equation}
				X\setminus K=\bigcup_{i\in \eN}(X\setminus K_i).
			\end{equation}
			Donc \( X\setminus K\) est ouvert en tant que union d'ouverts.


			\spitem[Un recouvrement pour \( K_0\)]
			%-----------------------------------------------------------
			Nous avons :
			\begin{subequations}
				\begin{align}
					K_0 & =(K_0\cap K)\cup(K_0\setminus K)                           \\
					    & \subset K\cup(K_0\setminus K)                              \\
					    & \subset\Big( \bigcup_{i\in I}A_i \Big)\cup (E\setminus K).
				\end{align}
			\end{subequations}
			Donc \( \{ A_i,E\setminus K \}\) est un recouvrement de \( K_0\) par des ouverts. Il existe donc un \( J\) fini dans \( \eN\) tel que
			\begin{equation}
				K_0\subset \big( \bigcup_{j\in J}A_j \big)\cup(X\setminus K).
			\end{equation}

			\spitem[Le recouvrement fini de \( K\)]
			%-----------------------------------------------------------
			Si \( x\in K\), alors \( x\in K_0\), mais \( x\) n'est pas dans \( X\setminus K\). Donc \( x\in \bigcup_{j\in J}A_j\). Nous avons démontré que \( \{ A_j \}_{j\in J}\) était un recouvrement fini de \( K\).
		\end{subproof}
	\end{subproof}
\end{proof}


\begin{lemma}[\cite{MonCerveau}]        \label{LEMooFJZDooSxYWVW}
	Toute union finie de compacts est compacte.
\end{lemma}

\begin{proof}
	Soient \( (K_i)_{i=1,\ldots, n}\) des compacts dans \( X\). Si \( \{ \mO_s \}_{i\in S}\) est un recouvrement de \( \bigcup_{i\in I}K_i\) par des ouverts, à fortiori, ce sera un recouvrement de chacun des \( K_i\). Pour chaque \( i\), il existera donc une partie finie \( S_i\) de \( S\) telle que \( \{ \mO_s \}_{s\in S_i}\) recouvre \( K_i\).

	L'union finie de parties finies \( S_i \) est une partie finie de \( S\), et nous avons
	\begin{equation}
		\bigcup_iK_i\subset \bigcup_{s\in \bigcup_iS_i}\mO_s.
	\end{equation}
\end{proof}

\begin{proposition}[\cite{BIBooRLJDooKCEKjj}]     \label{PROPooQWHSooXeJOkT}
	Dans un espace séparé, toute intersection de compacts est compacte.
\end{proposition}

\begin{proof}
	Soit un espace topologie séparé \( X\) et des compacts \( \{ K_i \}_{i\in I}\) dans \( X\) (\( I\) est un ensemble quelconque). Chacun des \( K_i\) est fermé par le lemme \ref{LemnAeACf}\ref{ITEMooAZWVooLyPDeY}. Donc l'intersection
	\begin{equation}
		K=\bigcap_{i\in I} K_i
	\end{equation}
	est un fermé de \( X\) par le lemme \ref{LemQYUJwPC}\ref{ITEMooBHIGooMvkUtX}. Soit \( i\) dans \( I\). Nous avons \( K\subset K_i\). Donc \( K\) est un fermé dans le compact \( K_i\); il est donc compact par le lemme \ref{LemnAeACf}.
\end{proof}

\begin{example}[Intersection de compacts non compacte\cite{BIBooRLJDooKCEKjj}]
	Un exemple d'intersection de compacts qui n'est pas compacte. Vu la proposition \ref{PROPooQWHSooXeJOkT}, il va falloir chercher un espace non séparé. Soit \( X=\eN\cup\{ x_1,x_2 \}\) où \( x_1\) et \( x_2\) sont deux éléments distincts hors de \( \eN\). Nous définissons une topologie sur \( X\) en disant que les ouverts sont les parties suivantes :
	\begin{itemize}
		\item les parties de \( \eN\),
		\item la partie \( \eN\cup\{ x_1 \}\),
		\item la partie \( \eN\cup\{ x_2 \}\),
		\item la partie \( \eN\cup\{ x_1,x_2 \}\).
	\end{itemize}
	Nous considérons les parties \( K_1=\eN\cup\{ x_1 \}\) et \(K_2= \eN\cup \{ x_2 \}\).
	\begin{subproof}
		\spitem[\( K_i\) est compact]
		Soit \( \{ \mO_i \}_{i\in I}\) un recouvrement de \( K_1\) par des ouverts de \( X\). Alors il existe \( i_0\in I\) tel que \( x_1\in\mO_{i_0}\). Vue la liste des ouverts, \( \mO_{i_0}\) est soit \( \eN\cup\{ x_1 \}\) soit \( \eN\cup\{ x_1,x_2 \}\). Dans les deux cas, \( \{ \mO_{i_0} \}\) est un sous-recouvrement fini de \( K_1\).
		\spitem[\( K_1\cap K_2=\eN\)]
		C'est immédiat parce que \( x_1\) et \( x_2\) sont distincts.
		\spitem[\( \eN\) n'est pas compact]
		Il peut être recouvert par les ouverts \( \{ \{ i \} \}_{i\in \eN}\) dont on ne peut pas extraire de sous-recouvrements finis.
	\end{subproof}
\end{example}

\begin{proposition}     \label{PropGBZUooRKaOxy}
	Si \( V\) est une partie de l'espace topologique \( X\) muni de la topologie induite\footnote{Définition \ref{DefVLrgWDB}.} \( \tau_V\) de celle de \( X\), et si \( K\) est un compact de \( (V,\tau_V)\) alors \( K\) est un compact de \( (X,\tau_X)\).
\end{proposition}

\begin{proof}
	Soient \( \{ \mO_\alpha \}_{\alpha\in A}  \) des ouverts de \( X\) recouvrant \( K\). Alors les ensembles \( V\cap \mO_{\alpha}\) recouvrent également \( K\), mais sont des ouverts de \( V\). Donc il en existe un sous-recouvrement fini. Soient donc \( \{V\cap\mO_i\}_{i\in I}\) recouvrant \( K\) avec \( I\) un sous-ensemble fini de \( A\). Les ensembles \( \{\mO_i\}_{i\in I}\) recouvrent encore \( K\) et sont des ouverts de \( X\).
\end{proof}

\begin{proposition}[\cite{MonCerveau}]
	Soient des espaces topologiques \( X\) et \( Y\). Nous considérons des ouverts \( A\) de \( X\) et \( B\) de \( Y\). Soit un compact \( M\) dans \( A\times B\)\footnote{Topologie produit, définition \ref{DefIINHooAAjTdY}}. Il existe des compacts \( K\) et \( L\) dans \( A\) et \( B\) tels que \( M\subset K\times L\).
\end{proposition}

\begin{proof}
	Nous considérons les «projections» de \( M\) sur \( A\) et \( B\):
	\begin{equation}
		K=\{ a\in A\tq \exists b\in B\tq (a,b)\in M \},
	\end{equation}
	et
	\begin{equation}
		L=\{ b\in B\tq \exists a\in A\tq (a,b)\in M \}.
	\end{equation}
	Nous avons \( M\subset K\times L\); il reste à montrer que \( K\) et \( L\) sont des compacts de leurs espaces respectifs. Soit un recouvrement \( \{ U_i \}_{i\in I} \) de \( K\) par des ouverts de \( X\) et \( \{ V_j \}_{j\in J}\) de \( L\) par des ouverts de \( Y\). Alors
	\begin{equation}
		\{ U_i\times  V_j \}_{\substack{i\in I\\j\in J}}
	\end{equation}
	est un recouvrement de \( M\) par des ouverts de \( X\times Y\). Puisque \( M\) est un compact de \( X\times Y\), nous pouvons en extraire un sous-recouvrement fini, c'est-à-dire \( I_0\) fini dans \( I\) et \( J_0\) fini dans \( J\) tels que
	\begin{equation}
		\{ U_i\times  V_j \}_{\substack{i\in I_0\\j\in J_0}}
	\end{equation}
	soit encore un recouvrement de \( K\times L\). Nous prouvons à présent que \( \{ U_i \}_{i\in I_0}\) est un recouvrement de \( K\), ce qui montrera que \( K\) est un compact.

	Soit \( a\in K\). Il existe \( b\in B\) tel que \( (a,b)\in M\). Donc il existe \( i_0\in I_0\) et \( j_0\in J_0\) tels que \( (a,b)\in U_{i_0}\times V_{j_0}\). En particulier \( a\in U_i\).

	Le même raisonnement montre que \( \{ V_j \}_{j\in J_0}\) est un recouvrement de \( L\).
\end{proof}


%---------------------------------------------------------------------------------------------------------------------------
\subsection{Espace relativement compact}
%---------------------------------------------------------------------------------------------------------------------------

\begin{definition}[Partie relativement compacte\cite{BIBooARERooBepIbf}]      \label{DEFooBODRooEFhzeT}
	Une partie d'un espace topologique est \defe{relativement compact}{relativement compact} si elle est contenue dans un compact.
\end{definition}

\begin{proposition}[\cite{BIBooARERooBepIbf}]	\label{PROPooXKBXooUIwTiO}
	Une partie d'un espace séparé est relativement compacte si et seulement si sa fermeture est compacte.
	%TODOooEOPQooYgjnZv. Prouver ça.
\end{proposition}

\begin{lemma}[Union finie de relativement compacts\cite{MonCerveau}]		\label{LEMooJXGUooXDGIXG}
	Dans un espace séparé\footnote{Séparé et Hausdorff sont synonymes, définition \ref{DefYFmfjjm}.}, une union finie de parties relativement compactes\footnote{Définition \ref{DEFooBODRooEFhzeT}.} est relativement compacte.
\end{lemma}

\begin{proof}
	Soit des paries relativement compactes \( A_i\). Notre objectif est de prouver que \( \overline{\bigcup_{i=1}^nA_i}\) est compact. La partie \( \bigcup_{i=1}^n\bar A_i\) est compacte comme union finie de compacts (lemme \ref{LEMooFJZDooSxYWVW}); elle est donc fermée par le lemme \ref{LemnAeACf}\ref{ITEMooAZWVooLyPDeY} :
	\begin{equation}
		\overline{\bigcup_{i=1}^n\bar A_i}=\bigcup_{i=1}^n\bar A_i.
	\end{equation}
	Nous avons donc
	\begin{equation}
		\overline{\bigcup_{i=1}^nA_i}\subset\overline{\bigcup_{i=1}^n\bar A_i}=\bigcup_{i=1}^n\bar A_i.
	\end{equation}
	Autrement dit, la partie \( \overline{\bigcup_{i=1}^nA_i}\) est un fermé dans un compact. Elle est donc compacte par le lemme \ref{LemnAeACf}\ref{ITEMooNKAKooQoNddr}.
\end{proof}

\begin{lemma}[Les ouverts séparent les points et les compacts\cite{MonCerveau}]			\label{LEMooNIJEooMdvkuA}
	Soit un espace séparé \( X\). Soient un compact \( K\) et un point \( a\not\in K\). Alors il existe des ouverts \( A\) et \( B\) tels que \( a\in A\), \( K\subset B\) et \( A\cap B=\emptyset\).
\end{lemma}

\begin{proof}
	Pour chaque \( x\in K\) nous séparons \( x\) et \( a\) : il existe \( A_x\) ouvert contenant \( a\), \( B_x\) ouvert contenant \( x\) tels que \( A_x\cap B_x=\emptyset\). Les \( \{ B_x \}_{x\in K}\) forment un recouvrement de \( K \) par des ouverts. Nous en extrayons un sous-recouvrement fini : il existe \( x_1,\ldots,x_n\in K\) tels que \( K\subset\bigcup_{i=1}^nB_{x_i}\). Nous notons \( A_i=A_{x_i}\) et \( B_i=B_{x_i}\).

	Nous posons
	\begin{equation}
		\begin{aligned}[]
			A=\bigcap_{i=1}^nA_i \\
			B=\bigcup_{i=1}^nB_i.
		\end{aligned}
	\end{equation}
	En tant qu'intersection finie d'ouverts, \( A\) est ouvert, et comme union d'ouverts, \( B\) est ouvert. Par construction, \( K\subset B\) et \( a\in A\). Il reste à voir que \( A\cap B=\emptyset\). Si \( y\in A\cap B\), nous avons \( y\in B\) et donc nous avons \( y\in B_j\) pour un certain \( j\). Mais \( y\in A_i\) pour tout \( i\), et en particulier \( y\in A_j\). Donc \( y\in A_i\cap B_i\). Impossible parce que nous avons choisit \( A_i\cap B_i=\emptyset\). Donc \( A\cap B\) est vide.
\end{proof}


%---------------------------------------------------------------------------------------------------------------------------
\subsection{Propriété d'intersection finie}
%---------------------------------------------------------------------------------------------------------------------------

\begin{definition}[Propriété d'intersection finie\cite{BIBooXACUooPPyLuN}]      \label{DEFooCESGooZkACqs}
	Soit un ensemble \( X\). Une famille non vide \( \mA\) de parties de \( X\) a la \defe{propriété d'intersection finie}{propriété d'intersection finie} si toutes intersection finie d'éléments de \( \mA\) est non vide.
\end{definition}

\begin{theorem}[\cite{BIBooXACUooPPyLuN,BIBooTYSSooWPqfAx}]       \label{THOooCQSQooDuasqo}
	Un espace est compact si et seulement si toute famille de parties fermées ayant la propriété d'intersection finie\footnote{Définition \ref{DEFooCESGooZkACqs}.} a une intersection non vide.
\end{theorem}

\begin{proof}
	En deux parties.
	\begin{subproof}
		\spitem[\( \Rightarrow\)]
		%-----------------------------------------------------------
		Nous supposons que \( X\) est un espace topologique compact et nous considérons une famille \( \{ F_i \}_{i\in I}\) de fermés ayant la propriété d'intersection finie. Nous supposons par l'absurde que \( \bigcap_{i\in I}F_i=\emptyset\) et nous allons contredire la propriété d'intersection finie.

		Nous avons, en utilisant le lemme \ref{LEMooMDTPooAhmYqv} :
		\begin{equation}
			X=\emptyset^x=\left( \bigcap_{i\in I}F_i \right)^c=\bigcup_{i\in I}F_i^c.
		\end{equation}
		Vu que les \( F_i\) sont fermés, les \( F_i^c\) sont ouverts. Donc \( \{ F_i^c \}_{i\in I}\) est un recouvrement de \( X\) par des ouverts. Par compacité, il existe \( J\) fini dans \( I\) tel que \( \bigcup_{j\in J}F_j^c=X\). Mais alors
		\begin{equation}
			X=\bigcup_{j\in J}F_j^c=\left( \bigcap_{j\in J}F_j \right)^c,
		\end{equation}
		de telle sorte que \( \bigcap_{j\in J}F_j=\emptyset\), qui contredit la propriété d'intersection finie.

		\spitem[\( \Leftarrow\)]
		%-----------------------------------------------------------
		Nous supposons que toute famille de fermés ayant la propriété d'intersection finie a une intersection finie. Et nous allons prouver que \( X\) est compact. Pour cela nous considérons un recouvrement \( \{ A_i \}_{i\in I}\) de \( X\) par des ouverts. Nous y allons par contradiction en disant que pour tout \( J\) fini dans \( I\) nous avons \( \bigcup_{j\in J}A_j\neq X\) (càs que \( X\) n'est pas compact), et nous allons trouver une contradiction.

		Si \( J\) est fini dans \( I\) nous avons
		\begin{equation}		\label{EQooCBAZooJoQamE}
			\bigcap_{j\in J}A_j^c=\left( \bigcup_{j\in J}A_j \right)^c\neq \emptyset.
		\end{equation}
		Mais les \( A_i^c\) sont des fermés, donc par propriété d'intersection finie,
		\begin{equation}
			\emptyset\neq\bigcap_{i\in I}A_i^c=\left( \bigcup_{i\in I}A_i \right)^c;
		\end{equation}
		la première inégalité est parce que la famille \( \{ A_i^c \}_{i\in I}\) sont des fermés a toutes ses intersections finies non vides (c'est \eqref{EQooCBAZooJoQamE}), et donc a une intersection non vide par hypothèse.

		Cela contredit le fait que \( \bigcup_{i\in I}A_i\) est un recouvrement de \( X\).
	\end{subproof}
\end{proof}

%+++++++++++++++++++++++++++++++++++++++++++++++++++++++++++++++++++++++++++++++++++++++++++++++++++++++++++++++++++++++++++
\section{Limite de fonction}
%+++++++++++++++++++++++++++++++++++++++++++++++++++++++++++++++++++++++++++++++++++++++++++++++++++++++++++++++++++++++++++

\begin{definition}[Limite d'une fonction, thème~\ref{THEMEooGVCCooHBrNNd}\cite{BIBooPGRGooBrqxAA}]\label{DefYNVoWBx}
	Soient des espaces topologiques \( X\) et \( Y\) ainsi que \( \Omega\subset X\) et \( a\in \Adh(\Omega)\). Soit une application \( f\colon \Omega\to Y\). Nous disons que l'élément \( \ell\) de \( Y\) est une \defe{limite}{limite!d'une fonction} de \( f\) en \( a\) lorsque pour tout ouvert \( V\) contenant \( \ell\), il existe un voisinage ouvert \( U\) de \( a\) tel que
	\begin{equation}        \label{EQooXLJJooZDcOtU}
		f\Big( \big( U\cap\Omega\big)\setminus\{ a \} \Big)\subset V.
	\end{equation}
	Si un tel élément est unique\footnote{Rappelons que ce n'est pas toujours le cas, mais que ça l'est si l'espace topologique est séparé -- définition~\ref{DefYFmfjjm}.}, alors nous disons que cet élément est la \defe{limite}{limite!d'une fonction} de \( f\) et nous notons
	\begin{equation}
		\lim_{x\to a} f(x)=\ell.
	\end{equation}
\end{definition}

\begin{normaltext}
	Il aurait été tout aussi bien de définir la limite d'une fonction \( f\colon X\to Y\) définie sur tout \( X\), puis de considérer \( \Omega\) avec la topologie induite depuis \( X\).

	Dans ce cas, nous aurions écrit \eqref{EQooXLJJooZDcOtU} sous la forme
	\begin{equation}
		f\Big(  U\setminus\{ a \} \Big)\subset V
	\end{equation}
	en disant que \( U\) est un voisinage de \( a\), et en laissant \randomGender{le lecteur}{la lectrice} deviner que ici, «voisinage» signifie «voisinage au sens de la topologie induite». Vu que nous considérons la fonction \( f\) uniquement définie sur \( \Omega\), c'est la seule interprétation possible, et il n'y aurait pas au d'ambigüité\cite{BIBooMDAKooGEtFUd}. Mais bon\ldots{} si ça va sans dire, ça va encore mieux en le disant.
\end{normaltext}

\begin{remark}
	Nous ne saurions trop insister sur le fait que la valeur de \( f\) en \( a\) n'intervient pas dans la définition de la limite de \( f\) en \( a\). Il n'est même pas nécessaire que \( f\) soit définie en \( a\) pour que l'on puisse parler de limite de \( f\) en \( a\). Par exemple nous avons
	\begin{equation}
		\lim_{x\to 1} \frac{ x^2-1 }{ x-1 }=2,
	\end{equation}
	alors que la fonction n'est pas définie en \( x=1\).

	Plus généralement, un peu par principe, toutes les fois que la notion de limite apporte une information, le point où l'on prend la limite est spécial. Sinon on ne calculerait pas la limite, mais on regarderait directement la valeur de la fonction. Cela est typiquement le cas lorsque nous aborderons les dérivées. En effet, regardons (en faisant semblant d'anticiper) la définition  \eqref{DEFooOYFZooFWmcAB}. Dans la formule
	\begin{equation}
		f'(a)=\lim_{x\to a} \frac{ f(x)-f(a) }{ x-a },
	\end{equation}
	la fonction sur laquelle nous prenons la limite n'est \emph{jamais} définie en \( x=a\).
\end{remark}

\begin{proposition}[Unicité de la limite pour un espace séparé]\label{PropFObayrf}
	Soient \( X\) un espace topologique, \( A\) une partie de \( X\) et \( Y\) un espace topologique séparé\footnote{Définition~\ref{DefYFmfjjm}.}. Nous considérons une fonction \( f\colon A\to Y\). Si \( a\in \Adh(A)\), alors \( f\) admet au plus une limite en \( a\).
\end{proposition}
\index{limite!unicité}

\begin{proof}
	Soient \( y\) et \( y'\) des limites de \( f\) en \( a\), ainsi que des voisinages \( V\) et \( V'\) de \( y\) et \( y'\). Nous prenons également les voisinages \( W\) et \( W'\) correspondants :
	\begin{subequations}
		\begin{numcases}{}
			f(W\cap A)  \subset V   \\
			f(W'\cap A) \subset V'.
		\end{numcases}
	\end{subequations}
	Quitte à prendre des sous-ensembles nous pouvons supposer que \( W\) et \( W'\) sont ouverts. Il s'ensuit alors que:
	\begin{itemize}
		\item l'ensemble \( W\cap W'\) est un ouvert contenant \( a\) et intersecte donc \( A\);
		\item l'ensemble \( (W\cap W')\cap A\) est donc non vide;
		\item et donc, \( f(W\cap W'\cap A) \) est, lui aussi, non vide.
	\end{itemize}
	Mais
	\begin{equation}
		f(W\cap W'\cap A)\subset f(W\cap A)\subset V,
	\end{equation}
	et
	\begin{equation}
		f(W\cap W'\cap A)\subset f(W'\cap A)\subset V',
	\end{equation}
	d'où \( V \) et \( V'\) ont une intersection. Puisque ces ensembles sont arbitraires, nous avons prouvé que tout voisinage de \( y\) et tout voisinage de \( y'\) ont une intersection non vide; étant donné que \( Y\) est séparé, nous devons avoir \( y=y'\).
\end{proof}

%+++++++++++++++++++++++++++++++++++++++++++++++++++++++++++++++++++++++++++++++++++++++++++++++++++++++++++++++++++++++++++
\section{Topologie, distances et normes}
%+++++++++++++++++++++++++++++++++++++++++++++++++++++++++++++++++++++++++++++++++++++++++++++++++++++++++++++++++++++++++++
Certains ensembles ont plus de structures qu'une topologie. Nous fixons quelques bases maintenant, et nous détaillerons certains résultats plus tard.

%---------------------------------------------------------------------------------------------------------------------------
\subsection{Distance et topologie métrique}
%---------------------------------------------------------------------------------------------------------------------------

\begin{definition}  \label{DefMVNVFsX}
	Si \( E\) est un ensemble, une \defe{distance}{distance} sur \( E\) est une application \( d\colon E\times E\to \eR\) telle que pour tout \( x,y\in E\),
	\begin{enumerate}

		\item
		      \( d(x,y)\geq 0\)

		\item
		      \( d(x,y)=0\) si et seulement si \( x=y\),

		\item
		      \( d(x,y)=d(y,x)\)

		\item
		      \( d(x,y)\leq d(x,z)+d(z,y)\).

	\end{enumerate}
	La dernière condition est l'\defe{inégalité triangulaire}{inégalité!triangulaire}.

	Un couple \( (E,d)\) formé d'un ensemble et d'une distance est un \defe{espace métrique}{espace!métrique}.
\end{definition}


\begin{definition}[Espace vectoriel topologique métrisable\cite{ooOFEPooVFgTXm}]		\label{DEFooGLTUooJjaSmI}
	Un espace topologique est \defe{métrisable}{métrisable!espace vectoriel topologique} si il existe une distance\footnote{Distance, définition \ref{DefMVNVFsX}.} compatible avec la topologie.
\end{definition}
\index{espace!vectoriel topologique!métrisable}


Le lemme suivant est similaire à la proposition \ref{PropNmNNm}.
\begin{lemma}[Inégalité triangulaire inversée\cite{BIBooMistral}]		\label{LEMooXCXHooVtrkvl}
	Si \( (E,d)\) est un espace métrique et si \( a,b,x\in E\), nous avons
	\begin{equation}
		d(a,b)\geq \big| d(a,x)-d(x,b) \big|
	\end{equation}
\end{lemma}

\begin{proof}
	Nous écrivons l'inégalité triangulaire pour \( d(a,x)\) et \( d(b,x)\) :
	\begin{subequations}
		\begin{numcases}{}
			d(a,x)\leq d(a,b)+d(b,x)\\
			d(b,x)\leq d(b,a)+d(a,x).
		\end{numcases}
	\end{subequations}
	Nous isolons alors \( d(a,b) \) dans chacune de ces inégalité, et nous utilisons la symétrie de \( d\) :
	\begin{subequations}
		\begin{numcases}{}
			d(a,b)\geq d(a,x)-d(x,b)\\
			d(a,b)\geq d(b,x)-d(a,x).
		\end{numcases}
	\end{subequations}
	En vertu du lemme \ref{LemooANTJooYxQZDw}\ref{ITEMooLWPTooEvHzvO}, nous avons alors
	\begin{equation}
		d(a,b)\geq |d(a,x)-d(x,b)|.
	\end{equation}
\end{proof}

La définition suivante donne une topologie sur les espaces métriques en partant des boules.

\begin{theoremDef}[Topologie métrique]     \label{ThoORdLYUu}
	Soit \( (E,d)\) un espace métrique. Nous définissons les \defe{boules ouvertes}{boule!ouverte} par
	\begin{equation}        \label{EQooYCWSooIhibvd}
		B(a,r)=\{ x\in E\tq d(a,x)<r \}.
	\end{equation}
	pour tout \( a\in E\) et \( r>0\).
	Alors en posant
	\begin{equation}        \label{EqGDVVooDZfwSf}
		\tau_d=\big\{  \mO\subset E  \tq\forall a\in \mO,\exists r>0\tq B(a,r)\subset \mO \big\}
	\end{equation}
	nous définissons une topologie sur \( E\).

	Cette topologie sur \( E\) est la \defe{topologie métrique}{topologie!métrique} de \( (E,d)\). En présence d'une distance, sauf mention explicite du contraire, c'est toujours cette topologie-là que nous utiliserons.
\end{theoremDef}

\begin{proof}
	D'abord \( \emptyset\in\tau_d\) parce que tout élément de l'ensemble vide \ldots{} heu \ldots{} enfin parce que, d'accord hein\footnote{Pour qui ne serait pas d'accord, ajoutez \( \emptyset\) dans la définition des ouverts et puis c'est tout.}. Ensuite si les \( \{A_i\}_{i\in I}\) sont des éléments de \( \tau_d\) et si \( x\in\bigcup_{i\in I}A_i\) alors il existe \( k\in I\) tel que \( x\in A_k\). Par hypothèse il existe une boule \( B(x,r)\subset A_k\subset\bigcup_{i\in I}A_i\).

	Enfin si les \( \{A_i\}_{i\in\{ 1,\ldots, n \}}\) sont des éléments de \( \mT\) alors pour tout \( i\) il existe \( r_i>0\) tel que \( B(x,r_i)\subset A_i\). En prenant \( r=\min\{ r_i \}_{i=1,\ldots, n}\) nous avons \( B(x,r)\subset\bigcap_{i=1}^nA_i.\)
\end{proof}


\begin{proposition}     \label{PROPooZXTXooEMLgMn}
	La topologie sur un espace métrique\footnote{Définition \ref{ThoORdLYUu}.} est la topologie engendrée\footnote{Définition \ref{DefTopologieEngendree}} par ses boules ouvertes.
\end{proposition}

\begin{proof}
	Nous notons \( (E,d)\) l'espace métrique. Nous notons \( \tau_d\) la topologie métrique, et \( \tau_B\) celle engendrée par les boules ouvertes.
	\begin{subproof}
		\spitem[\( \tau_d\subset\tau_B\)]
		%-----------------------------------------------------------
		Soit \( a\in\tau_d\). Pour tout \( a\in A\), nous avons une boule \( B(a,r_a)\) inclue dans \( A\), et donc
		\begin{equation}
			A=\bigcup_{a\in A}B(a,r_a)\in \tau_B.
		\end{equation}

		\spitem[\( \tau_B\subset \tau_d\)]
		%-----------------------------------------------------------
		Soient \( a_i\in E\) et \( r_i>0\).

		\begin{subproof}
			\spitem[\( \bigcap_{i=1}^nB(a_i,r_i)\in \tau_d\)]
			%-----------------------------------------------------------
			Soit \( x\in\bigcap_{i=1}^nB(a_i,r_i)\). Pour chaque \( i\), nous avons \( d(x,a_i)<r_i\). Il existe donc \( \alpha_i>0\) tel que \( d(x,a_i)<\alpha_i<r_i\). En posant \( \epsilon_i<r_i-\alpha_i\) nous avons \( B(x,\epsilon_i)\subset B(a_i,r_i)\). En posant \( \epsilon<\min_i(\epsilon_i)\) nous avons
			\begin{equation}
				B(x,\epsilon)\subset\cap_{i=1}^nB(a_i,r_i).
			\end{equation}
			\spitem[Et donc]
			%-----------------------------------------------------------
			Nous avons prouvé que
			\begin{equation}
				\bigcap_{i=1}^nB(a_i,r_i)\in\tau_d.
			\end{equation}
			Tout élément de \( \tau_B\) est une union de tels éléments. Vu que \( \tau_d\) est stable par union, nous avons \( \tau_B\subset\tau_d\).
		\end{subproof}
	\end{subproof}
\end{proof}

\begin{proposition}
	À propos de séparation.
	\begin{enumerate}
		\item
		      Tout espace métrique\footnote{Définition \ref{DefMVNVFsX}.} est séparé.
		\item
		      Si une suite dans un espace métrique possède une limite, alors elle est unique.
	\end{enumerate}
\end{proposition}

\begin{proof}
	Si deux éléments \( x \) et \( y \) sont distincts, alors en posant \( r = d(x , y) / 3 > 0 \), les boules \( B(x,r) \) et \( B(y,r)\) sont disjointes.

	En ce qui concerne les limites, ce sont les propositions \ref{PropUniciteLimitePourSuites} et \ref{PropFObayrf}.
\end{proof}

\begin{normaltext}
	Nous allons maintenant voir deux résultats disant que si une fonction est continue, alors elle peut être permutée avec une limite de suite. Dans le cas des espaces métriques, la proposition \ref{PropXIAQSXr} montrera la réciproque : si pour toute suite \(x_n\to a\), nous avons \( \lim_{n\to \infty} f(x_n)=y\), alors \( f\) a une limite en \( a\) qui vaut \( y\).
\end{normaltext}

\begin{proposition}[Permuter limite et fonction continue\cite{MonCerveau}] \label{fContEstSeqCont}
	Soient deux espaces topologiques \( X\) et \( Y\) ainsi qu'une fonction \( f\colon X\to Y\). Soit \( a\in X\) et \( \ell\in Y\). Si
	\begin{equation}
		\lim_{x\to a} f(x)=\ell,
	\end{equation}
	alors, pour toute suite \( (x_k) \) telle que \( x_k \to a \), on a
	\begin{equation}
		\lim f(x_k)=\ell.
	\end{equation}
\end{proposition}

\begin{proof}
	Nous considérons une suite \( (x_k)\) qui converge vers \( a\) dans \( X\). Soient \( V\) un voisinage de \( \ell \) et \( W\) un voisinage de \( a\) tels que \( f(W)\subset V\) (définition~\ref{DefYNVoWBx} de la continuité en un point). Par la convergence \( x_k\to a\),  il existe \( N\) tel que pour tout \( k\geq N\), \( x_k\in W\), et donc tel que \( f(x_k)\in V\), ce qui donne la continuité séquentielle de \( f\).
\end{proof}


%-------------------------------------------------------
\subsection{Précompact}
%----------------------------------------------------


\begin{definition}	\label{DEFooLZDTooZlAtdL}
	Une partie \( A\) d'un espace topologique métrique\footnote{Topologie métrique, définition \ref{ThoORdLYUu}.} est \defe{précompact}{précompact} si pour tout \( \epsilon>0\), il peut être recouvert par un nombre fini de boules de rayon \( \epsilon\).
\end{definition}

\begin{proposition}[\cite{BIBooKKSLooGjoYjh}]	\label{PROPooDWUDooGrgZpP}
	L'adhérence d'un précompact est précompacte.
\end{proposition}

\begin{proof}
	Soit un précompact \( A\); fixons \( \epsilon>0\). La partie \( A\) peut être recouverte par les boules \( \{ B(a_i,\epsilon) \}_{i=1,\ldots,n}\). Soit maintenant \( x\in\Adh(A)\). Il existe \( a\in A\) tel que \( d(a,x)<\epsilon\). De plus ce \( a\) est dans une des boules; c'est-à-dire que il existe \( i=1,\ldots,n\) tel que \( d(a_i,a)<\epsilon\). Nous avons donc
	\begin{equation}
		d(a_i,x)\leq d(a_i,a)+d(a,x)\leq 2\epsilon.
	\end{equation}
	Donc les boules \( \{ B(a_i,2\epsilon) \}_{i=1,\ldots,n}\) recouvrent \( \Adh(A)\).
\end{proof}

\begin{proposition}	\label{PROPooZRXDooUyFBFG}
	Soient un espace topologique \( X\) ainsi que \( A\subset X\). Si \( A\) est précompact et complet, alors \( A\) est compact.
	%TODOooGKFUooBmyebj. Prouver ça.
\end{proposition}



%-------------------------------------------------------
\subsection{Caractérisations séquentielles}
%----------------------------------------------------


\subsubsection{Continuité séquentielle}
%///////////////////////

\begin{definition}  \label{DefENioICV}
	Si \( X\) et \( Y \) sont deux espaces topologiques, une fonction \( f\colon X\to \eR\) est \defe{séquentiellement continue}{continuité!séquentielle} en un point \( a\) si pour toute suite convergente \( x_n\to a\) dans \( X\) nous avons \( f(x_n)\to f(a)\) dans \( Y\).
\end{definition}




%---------------------------------------------------------------------------------------------------------------------------
\subsection{Suites et espaces métriques}
%---------------------------------------------------------------------------------------------------------------------------

%TODO : il y a un contre-exemple à faire à la page http://www.les-mathematiques.net/phorum/read.php?14,787368,787582

\begin{proposition}[Caractérisation séquentielle de la limite\cite{MonCerveau}]     \label{PROPooJYOOooZWocoq}
	Soient deux espaces métriques \( X\) et \( Y\) ainsi qu'une fonction \( f\colon X\to Y\). Soit \( a\in X\) et \( \ell\in Y\). On a
	\begin{equation}\label{EqLimooJYOOooZWocoqG}
		\lim_{x\to a} f(x)=\ell,
	\end{equation}
	si et seulement si, pour toute suite \( (x_k) \) telle que \( x_k \to a \), on a
	\begin{equation}\label{EqLimooJYOOooZWocoqS}
		\lim f(x_k)=\ell.
	\end{equation}
	Par ailleurs, l'une des deux limites existe si et seulement si l'autre existe.
\end{proposition}

\begin{proof}
	Le sens direct est la proposition~\ref{fContEstSeqCont}. Pour la réciproque, nous passons par la contraposée. C'est-à-dire que nous supposons que \( \ell\) n'est pas une limite de \( f\) pour \( x\to a\). Il existe un \( \epsilon\) tel que pour tout \( \delta\), il existe un \( x\) vérifiant \( d_X(x;a) <\delta\) et \( d_Y(f(x);\ell) >\epsilon\).

	Nous construisons à présent une suite de la manière suivante. Pour \( \delta=\frac{1}{ n }\) nous considérons \( x_n\) tel que \( d_X( x_n; a) <\delta\) et \( d_Y(f(x_n);\ell) > \epsilon \). Cette suite converge vers \( a\), mais la suite \( f(x_n)\) ne converge manifestement pas vers \( \ell\) : elle ne rentre jamais dans la boule \( B(\ell,\epsilon)\).
\end{proof}

Une fonction continue est séquentiellement continue. Dans les espaces métriques la proposition suivante montre que la réciproque est également vraie et la continuité est équivalente à la continuité séquentielle. Cela n'est cependant pas vrai pour n'importe quel espace topologique.

\begin{corollary}[Caractérisation séquentielle de la continuité en un point\cite{MonCerveau}]  \label{PROPooBHRBooJMZYSg}
	Si \( X\) et \( Y\) sont des espaces métriques, alors une fonction \( f\colon X\to Y\) est continue en un point si et seulement si elle est séquentiellement continue en ce point.
\end{corollary}

\begin{proof}
	Paraphrasons la preuve précédente. Nous supposons que \( X\) et \( Y\) sont métriques. Si \( f\) n'est pas continue en \( a\), il existe \( \epsilon>0\) tel que pour tout \( \delta>0\), il existe \( x\) tel que \( \| x-a \|\leq\delta\) et \( \| f(x)-f(a) \|>\epsilon\). Nous considérons un tel \( \epsilon\) et pour chaque \( n\geq1\in \eN\) nous considérons un \( x_n\) correspondant à \( \delta=\frac{1}{ n }\). Cela nous donne une suite \( x_n\to a\) dans \( X\) mais \( \| f(x_n) -f(a)\|\) reste plus grand que \( \epsilon\). Cela montre que \( f\) n'est pas non plus séquentiellement continue.
\end{proof}


\begin{definition}[Fermeture séquentielle\cite{BIBooWWPDooIgMrch}]
	Une partie \( F\) d'un espace topologique \( X\) est dit \defe{séquentiellement fermé}{séquentiellement fermé} si la convergence d'une suite \( (x_n)\) de \( F\) vers \( x\) implique que \( x\) appartient à \( F\).
\end{definition}

Les espaces métriques ont une propriété importante que la fermeture séquentielle est équivalente à la fermeture.

\begin{proposition}[Caractérisation séquentielle d'un fermé]    \label{PropLFBXIjt}
	Soient \( X\) un espace métrique\footnote{Définition \ref{DefMVNVFsX}.} et \( F\subset X\). L'ensemble \( F\) est fermé si et seulement si toute suite contenue dans \( F\) et convergeant dans \( X\) converge vers un élément de \( F\).
\end{proposition}
\index{fermeture séquentielle}
\index{séquentiellement fermé}

\begin{proof}
	Une suite contenue dans un fermé ne peut converger que vers un élément de ce fermé: c'était la proposition \ref{PROPooBBNSooCjrtRb}. Le point le plus important est donc l'autre sens: si toute suite d'éléments de \( F \) converge dans \( F \) alors \( F \) est fermé.

	Par contraposée, supposons que \( X\setminus F\) ne soit pas ouvert. Alors il existe \( x\in X\setminus F\) pour lequel tout voisinage intersecte \( F\). En prenant \( x_k\in B(x,\frac{1}{ k })\), nous construisons une suite contenue dans \( F\), convergeant vers \( x\) qui n'est pas dans \( F \).
\end{proof}


Le lemme suivant est précisément la version «espace métrique» du corolaire \ref{CorLimAbarA}; mais, donnons-en une preuve tout de même.
\begin{lemma}		\label{LemLimAbarA}
	Soit \( X\) un espace métrique\footnote{Définition \ref{DefMVNVFsX}.}, et soit \( (x_n)\) une suite convergente contenue dans un ensemble \( A\subset X\). Alors la limite \( x_n\) appartient à \( \bar A\).
\end{lemma}

\begin{proof}
	Supposons que nous ayons une partie \( A\) de \( X\), et une suite \( (x_n)\) dont la limite \( \ell\) se trouve hors de \( \bar A\). Dans ce cas, il existe un \( r>0\) tel que\footnote{Une autre manière de dire la même chose : si \( \ell\notin\bar A\), alors \( d(\ell,A)>0\).} \( B(\ell,r)\cap A=\emptyset\). Si tous les éléments \( x_n\) de la suite sont dans \( A\), il n'y en a donc aucun tel que \( d(x_n,\ell)<r\). Cela contredit la notion de convergence \( x_n\to \ell\).
\end{proof}

\begin{corollary}		\label{CorAdhEstLim}
	Soit \( X\) un espace métrique, \( A \subset X\) et \( a \in \bar A\). Alors il existe une suite d'éléments dans \( A\) qui converge vers \( a\).
\end{corollary}

\begin{proof}
	Si \( a\in A\), alors nous pouvons prendre la suite constante \( x_n=a\). Si \( a\) n'est pas dans \( A\), alors \( a\) est dans \( \partial A\), et pour tout \( n\), il existe un point de \( A\) dans la boule \( B(a,\frac{1}{ n })\). Si nous nommons \( x_n\) ce point, la suite ainsi construite est une suite contenue dans \( A\) et qui converge vers \( a\) (ce dernier point est laissé à la sagacité \randomGender{du lecteur}{de la lectrice}.).
\end{proof}

En termes savants, ce corolaire signifie que la fermeture \( \bar A\) est composé de \( A\) plus de toutes les limites de toutes les suites contenues dans \( A\).

\begin{proposition}[Caractérisation séquentielle de la continuité\cite{MonCerveau}]     \label{PropXIAQSXr}
	Soient \( X\) et \( Y\) deux espaces topologiques séparés. Nous supposons que \( X\) est métrisable\footnote{Définition \ref{DEFooGLTUooJjaSmI} : il existe une distance compatible avec la topologie.}. Une application \( f\colon X\to Y\) est continue sur \( X\) si et seulement si elle est séquentiellement continue sur \( X\).
\end{proposition}

\begin{proof}
	En deux parties dont une déjà faite.
	\begin{subproof}
		\spitem[\( \Rightarrow\)]
		%-----------------------------------------------------------
		C'est la proposition \ref{fContEstSeqCont}.

		\spitem[\( \Leftarrow\)]
		%-----------------------------------------------------------
		Soit \( \mO\) un ouvert de \( Y\); nous allons voir que le complémentaire de \( f^{-1}(\mO)\) est fermé dans \( X\). Pour cela nous considérons une suite convergente \( x_k\stackrel{X}{\longrightarrow} x\) avec \( x_k\in X\setminus f^{-1}(\mO)\) pour tout \( k\). Nous allons montrer que \( x\in X\setminus f^{-1}(\mO)\) et la caractérisation séquentielle\footnote{Proposition~\ref{PropLFBXIjt}, valable parce que la topologie de \( X\) provient d'une métrique.} de la fermeture conclura que \( X\setminus f^{-1}(\mO)\) est fermé.

		Pour tout \( k\), nous avons \( f(x_k)\in X\setminus \mO\), et \( f(x_k)\stackrel{Y}{\longrightarrow} f(x)\) parce que \( f\) est séquentiellement continue. Vu que \( f(x_k)\) est une suite dans le fermé \( Y\setminus \mO\), la limite est également dans \( Y\setminus \mO\). Nous en déduisons que \( f(x)\in Y\setminus \mO\), de telle sorte que \( x\in X\setminus f^{-1}(\mO)\).
	\end{subproof}
\end{proof}

\begin{proposition} \label{PropCJGIooZNpnGF}
	Si \( X\) et \( Y\) sont deux espaces métriques et \( f,g\colon X\to Y\) sont deux fonctions continues égales sur une partie dense de \( X\) alors \( f=g\).
\end{proposition}
\index{fonction!continue!égales}

\begin{proof}
	Les fonctions \( f\) et \( g\) sont séquentiellement continues (proposition \ref{PROPooBHRBooJMZYSg}). Soient \( A\) un ensemble dense dans \( X\) sur lequel \( f\) et \( g\) sont égales, et \( x\notin A\). Vu que \( A\) est dense, il existe une suite \( a_n\) dans \( A\) telle que \( a_n\to x\). La séquentielle continuité de \( f\) et \( g\) donnent
	\begin{subequations}
		\begin{align}
			f(a_n)\to f(x) \\
			g(a_n)\to g(x),
		\end{align}
	\end{subequations}
	mais pour tout \( n\), \( f(a_n)=g(a_n)\). Par unicité de la limite\footnote{Proposition~\ref{PropFObayrf}.} dans \( Y\), \( f(x)=g(x)\).
\end{proof}



%-------------------------------------------------------
\subsection{Continuité de la distance}
%----------------------------------------------------


\begin{proposition}[\cite{MonCerveau}]	\label{PROPooSZMLooNdtLLj}
	Soit un espace métrique \( (E,d)\). L'application \(d \colon E\times E\to E  \) est continue.
	%TODOooZHSJooBbCibF.
\end{proposition}


\begin{proposition}[\cite{MonCerveau}]      \label{PROPooUXDJooCrWBbd}
	Soient un espace métrique \( (E,d)\), ainsi qu'une suite convergente \( a_n\stackrel{d}{\longrightarrow}\ell\). Il existe \( r>0\) tel que pour tout \( n\) nous ayons \( d(\ell, a_n)<r\).
\end{proposition}

\begin{proof}
	Soit \( r_1>0\) et \( N\in \eN\) tel que \( n\geq N\) implique \( d(\ell,a_n)<r_1\). Ensuite nous posons \( r_2=\max\{ d(\ell,a_n) \}_{n=0,\ldots, N}\).

	Pour tout \( n\) nous avons \( d(a_n,\ell)\leq r_1+r_2\).
\end{proof}

\begin{proposition}[\cite{MonCerveau}]		\label{PROPooNOHQooTqBQLk}
	Soient un espace topologique \( X\), un espace topologique normé\footnote{Topologie en \ref{ThoORdLYUu}.} \( Y\). Soit une application continue \(f \colon X\to Y  \). Pour tout \( x\in X\), il existe un voisinage ouvert de \( x\) sur lequel \( f\) est bornée.
\end{proposition}

\begin{proof}
	Soit \( \delta>0\). Nous considérons dans \( Y\) la boule \( B\big( f(x),\delta \big)\). Par continuité de \( f\), la partie \( f^{-1}\Big( B\big( f(x),\delta \big) \Big)\) est ouverte dans \( X\). Si \( y\) est dans cette partie, alors
	\begin{equation}
		f(y)\in B\big( f(x),\delta \big).
	\end{equation}
	L'application \( f\) y est donc bornée.
\end{proof}

%---------------------------------------------------------------------------------------------------------------------------
\subsection{Topologie métrique et induite}
%---------------------------------------------------------------------------------------------------------------------------

\begin{lemma}       \label{LEMooKDMYooMIcFRI}
	Soit un espace vectoriel normé \( \big( V,\| . \| \big)   \) muni d'un sous-espace vectoriel \( M\). La topologie induite de \( M\) depuis \( V\) est la même que la topologie de la norme induite \( \big( M,\| . \| \big)\).
	%TODOooDEAIooRhIKIr. Prouver ça.
\end{lemma}

%---------------------------------------------------------------------------------------------------------------------------
\subsection{Intérieur, adhérence et frontière}
%---------------------------------------------------------------------------------------------------------------------------

\begin{normaltext}
	Choses déjà faites :
	\begin{itemize}
		\item
		      Intérieur, définition \ref{DEFooSVWMooLpAVZRInt}.
		\item
		      Adhérence, qui est la même chose que fermeture, définition \ref{DEFooSVWMooLpAVZR}, et précisé par le lemme~\ref{LEMooILNCooOFZaTe}.
	\end{itemize}

	Dans le cas de \( \eR^n\) dans lequel les boules forment une base de la topologie nous pouvons encore préciser de la façon suivante:
	\begin{equation}
		x \in \Adh A \iffdefn \forall \epsilon > 0, B(x,\epsilon) \cap A \neq \emptyset
	\end{equation}
\end{normaltext}

\begin{proposition}		\label{PROPooZRFZooVqDFWD}
	Si \( A\) est une partie d'un espace topologique nous avons:
	\begin{equation}
		\Int A \subseteq A  \subseteq \Adh(A)
	\end{equation}
\end{proposition}

\begin{proof}
	Si \( x\in \Int(A)\), alors \( x\) admet un voisinage contenu dans \( A\); ledit voisinage contient \( x\), et donc \( x\) est dans \( A\).

	Si \( x\in A\) alors tout voisinage de \( x\) intersecte \( A\) (au moins en \( x\) lui-même), donc \( x\in\Adh(A)\).
\end{proof}

\begin{definition}      \label{DEFooACVLooRwehTl}
	La \defe{frontière}{frontière} ou le \defe{bord}{bord} de \( A\) est défini par \( \partial A = \Adh A \setminus \Int A\).
\end{definition}

\begin{lemma}       \label{LEMooMPZWooGrqYIX}
	Soit une partie \( A\) d'un espace topologique.
	\begin{enumerate}
		\item		\label{ITEMooNBSCooXJcVos}
		      \( A\) est ouverte si et seulement si \( A = \Int(A)\).
		\item		\label{ITEMooDCIOooQAvCqI}
		      \( A\) est fermée si et seulement si \( A = \Adh(A)\).
	\end{enumerate}
	Et je rappelle que «adhérence» et «fermeture» c'est la même chose et que \( \bar A\) et \( \Adh(A)\) sont deux notations pour la fermeture.
\end{lemma}

\begin{proof}
	Soient un espace topologique \( X\) et \( A\subset X\).
	\begin{subproof}
		\spitem[Pour \ref{ITEMooNBSCooXJcVos}]
		%-----------------------------------------------------------
		C'est une reformulation du théorème \ref{ThoPartieOUvpartouv} parce que \( x\) a un voisinage dans \( A\) signifie exactement \( x\in \Int(A)\).

		\spitem[Pour \ref{ITEMooDCIOooQAvCqI}]
		%-----------------------------------------------------------
		En deux parties.
		\begin{subproof}
			\spitem[Si \( A\) est fermé]
			%-----------------------------------------------------------
			Nous devons prouver que \( A=\bar A\) en sachant que \( A^c\) est ouvert. Nous écrivons la proposition \ref{PROPooTTWBooOyxWzr} avec \( A^c\) :
			\begin{equation}
				X=\Int(A^c)\cup \overline{X\setminus A^c}=A^c\cup \bar A.
			\end{equation}
			Nous avons utilisé le fait que \( \Int(A^c)=A^c\) en vertu de \ref{ITEMooNBSCooXJcVos} et que \( A^c\) est ouvert. Étant donné que l'union est disjointe, \( A^c=(\bar A)^c\), ce qui signifie que \( A=\bar A\) comme il fallait.

			\spitem[Si \( A=\bar A\)]
			%-----------------------------------------------------------
			Nous écrivons encore \ref{PROPooTTWBooOyxWzr} avec \( A^c\) :
			\begin{equation}
				X=\Int(A^c)\cup \bar A=\Int(A^c)\cup A.
			\end{equation}
			Vu que l'union est disjointe, \( \Int(A^c)=A^c\), ce qui signifie que \( A\) est ouvert par \ref{ITEMooNBSCooXJcVos}.
		\end{subproof}
	\end{subproof}
\end{proof}

\begin{proposition}[\cite{MonCerveau}]	\label{PROPooFHURooGXcekV}
	Un ouvert n'intersecte pas sa frontière.
\end{proposition}

\begin{proof}
	Soit un ouvert \( A\). Nous avons \( \partial(A)=\bar A\setminus\Int(A)\). Oh mais comme \( A\) est ouvert, \(  \Int(A)=A\). Donc les éléments de \( A\) sont exclus de \( \partial A\).
\end{proof}

\begin{proposition}[\cite{MonCerveau}]	\label{PROPooGXYCooRcnHDD}
	Si \( A\) est une partie de l'espace topologique \( X\), alors
	\begin{enumerate}
		\item
		      \begin{equation}
			      X\setminus \Int(A)=\overline{X\setminus A}
		      \end{equation}
		\item
		      \begin{equation}		\label{EQooOKINooXjFiUX}
			      \partial A=\bar A\setminus\overline{X\setminus A}.
		      \end{equation}
	\end{enumerate}
\end{proposition}

\begin{proof}
	Soit \( x\in X\setminus\Int(A)\). Vu que \( x\) n'est pas dans l'intérieur de \( A\), tout voisinage de \( x\) contient un point hors de \( A\), c'est à dire un point de \( X\setminus A\). Donc \( x\in\overline{X\setminus A}\).

	Dans l'autre sens si \( x\in\overline{X\setminus A}\), tout voisinage de \( x\) contient un élément de \( X\setminus A\). Donc \( x\) n'est pas dans \( \Int(A)\).

	En ce qui concerne \eqref{EQooOKINooXjFiUX}. Il suffit d'utiliser ce qu'on vient de faire :
	\begin{subequations}
		\begin{align}
			\partial A & =\bar A\setminus \Int(A)                  \\
			           & =\bar A\cap\big( X\setminus \Int(A) \big) \\
			           & =\bar A\cap\overline{X\setminus A}
		\end{align}
	\end{subequations}
\end{proof}

\begin{proposition}[\cite{MonCerveau}]	\label{PROPooTHOPooCOGmZD}
	Toute frontière est fermée. Plus précisément : si \( S\) est une partie de l'espace topologique \( X\), nous avons :
	\begin{enumerate}
		\item		\label{ITEMooCIGTooOyNTYt}
		      \( \partial S=\bar S\cap \overline{E\setminus S}\).
		\item		\label{ITEMooCSHKooHibSre}
		      \( \partial S\) est fermée.
	\end{enumerate}
\end{proposition}


\begin{proof}
	Pour \ref{ITEMooCIGTooOyNTYt}, nous utilisons la proposition \ref{PROPooTTWBooOyxWzr} : \( \overline{X\setminus S}=X\setminus\Int(S)\). Donc
	\begin{equation}
		\overline{S}\cap\overline{X\setminus S}=\bar S\cap\big( X\setminus \Int(S) \big)=\bar S\setminus \Int(S)=\partial(S).
	\end{equation}
	Nous avons utilisé le lemme \ref{LEMooHRKAooRskzQL}\ref{ITEMooRXJOooUYPaCp}.

	En ce qui concerne \ref{ITEMooCSHKooHibSre}, \( \partial S\) est fermé en tant qu'intersection de deux fermés.
\end{proof}


\begin{proposition}[frontière de l'union et de l'intersection\cite{BIBooKXTDooOmBoAj}]	\label{PROPooKQXDooEnNJyn}
	Soient des parties \( A\) et \( B\) d'un espace topologique \( X\). Nous avons
	\begin{equation}
		\partial(A\cup B)\subset \partial A\cup\partial B
	\end{equation}
	et
	\begin{equation}
		\partial(A\cap B)\subset \partial A\cup\partial B.
	\end{equation}
\end{proposition}

\begin{proof}
	Voici le calcul pour l'union :
	\begin{subequations}
		\begin{align}
			\partial(A\cup B) & =\overline{A\cup B}\cap\big( \overline{X\setminus(A\cup B)} \big)                                                            & \text{prop. \ref{PROPooGXYCooRcnHDD}.} \\
			                  & =\overline{A\cup B}\cap\big( \overline{(X\setminus A)\cap(X\setminus B)} \big)                                                                                        \\
			                  & \subset\overline{A\cup B}\cap\big( \overline{X\setminus A}\cap\overline{X\setminus B} \big)                                  & \text{prop. \ref{PROPooEQYEooZHXKnL}.} \\
			                  & =\bar A\cup\bar B\cap\overline{(X\setminus A)}\cap\overline{(X\setminus B)}                                                                                           \\
			                  & =\big( \bar A\cap\overline{X\setminus A} \big)\cup\big( \bar B\cap\overline{X\setminus A} \big)\cap\overline{X\setminus B}                                            \\
			                  & =\big(\partial A\cap\overline{X\setminus B}\big)\cup\big( \bar B\cap\overline{X\setminus A}\cap\overline{X\setminus B} \big) & \text{prop. \ref{PROPooGXYCooRcnHDD}}  \\
			                  & =\big( \partial A\cap\overline{X\setminus B} \big)\cup\big( \partial B\cap \overline{X\setminus A}\big)                                                               \\
			                  & \subset\partial A\cup\partial B.
		\end{align}
	\end{subequations}

	Et le calcul pour l'intersection :
	\begin{subequations}
		\begin{align}
			\partial(A\cap B) & =\overline{A\cap B}\cap\overline{\big(X\setminus(A\cap B)\big)}                                                                                               \\
			                  & =\overline{A\cap B}\cap \Big( \overline{(X\setminus A)\cup(X\setminus B)} \Big)                                                                               \\
			                  & =\overline{A\cap B}\cap\big( \overline{(X\setminus A)}\cup \overline{(X\setminus B)} \big)                              & \text{par \eqref{EQooSBGTooGPPbCl}} \\
			                  & =\big( \overline{A\cap B}\cap\overline{X\setminus A} \big)\cup\big( \overline{A\cap B}\cap\overline{X\setminus B} \big)                                       \\
			                  & \subset \big( \bar A\cap\overline{X\setminus A} \big)\cup \big( \bar B\cap\overline{X\setminus B} \big)                                                       \\
			                  & =\partial A\cup\partial B.
		\end{align}
	\end{subequations}
\end{proof}


\begin{proposition}[\cite{MonCerveau}]	\label{PROPooIIRLooIPdjqO}
	Soient un espace topologique \( X\) ainsi que \( A,B\subset X\). Soit \( x\in\partial A\). Nous suppposons qu'il existe un voisinage ouvert \( V_0\) de \( x\) tel que \( V_0\cap B=\emptyset\). Alors \( x\in\partial(A\cup B)\).
\end{proposition}

\begin{proof}
	Vu que \( x\in\partial A\), nous avons \( x\in\bar A\subset \overline{A\cup B}\). Il nous reste à prouver que \( x\) n'est pas dans l'intérieur de \( A\cup B\). Soit un voisinage ouvert \( W\) de \( x\). La partie \( W\cap V_0\) est un ouvert contenant \( x\) et tel que \( W\cap V_0=\emptyset\). De plus, étant donné que \( x\in\partial A\), il existe \( s\in W\cap V_0\) tel que \( s\not\in A\). Donc \( s\not\in A\cup B\). Et par conséquent il n'existe pas de voisinage de \( x\) contenu dans \( A\cup B\); autrement dit, \( x\) n'est pas dans l'intérieur de \( A\cup B\).
\end{proof}

\begin{proposition}[\cite{MonCerveau}]	\label{PROPooYHIGooMDVKNy}
	Soit une partie \( S\) d'un espace topologique \( X\). Soit un ouvert \( A\subset S\). Si \( A\cap\partial S=\emptyset\), alors \( \bar A\subset S\).
	%TODOooHLVYooWMQEoT. Prouver ça.
\end{proposition}



\begin{lemma}[Caractérisation équivalente de la frontière]      \label{LEMooEUYEooYcUfKr}
	Soient \( X\) un espace topologique et \( S\subset X\). Un point \( x\in X\) est dans \( \partial S\) si et seulement si tout voisinage de \( x\) contient un point de \( S\) et un point de \( S^c\).
\end{lemma}

\begin{proof}
	Supposons que tout voisinage de \( x\) contienne un point de \( S\) et un point de \( S^c\). Alors \( x\in Adh(S)\) (définition~\ref{DEFooSVWMooLpAVZR}), mais pas dans l'intérieur de \( S\) parce que \( x\) ne possède pas de voisinage contenu dans \( S\). Donc \( x\in \partial S\).

	À l'inverse, si \( x\in\partial S\) alors \( x\) est dans l'adhérence de \( S\) et tout voisinage de \( x\) contient un point de \( S\). Mais \( x\) n'est pas dans l'intérieur de \( S\) et tout voisinage de \( x\) contient un point qui n'est pas dans \( S\), aka un point de \( S^c\).
\end{proof}

\begin{corollary}
	Un ensemble et son complémentaire ont même frontière.
\end{corollary}

\begin{proof}
	Conséquence du lemme~\ref{LEMooEUYEooYcUfKr}. Les points de \( \partial(S^c)\) sont caractérisés par le fait que tout voisinage contient un point de \( S^c\) et un point de \( (S^c)^c=S\).
\end{proof}

\begin{example}
	Soit \( X=\mathopen[ 0 , 1 \mathclose]\) muni de la topologie de la distance \( | x-y |\) (définition~\ref{ThoORdLYUu}). Les points \( 0\) et \( 1\) \emph{ne sont pas} dans la frontière de \( X\). En effet une boule ouverte autour de \( 1\) est un ensemble de la forme
	\begin{equation}
		B(1,r)=\{ x\in X\tq | x-1 |<r \}=\mathopen] 1-r , 1 \mathclose]
	\end{equation}
	où nous avons supposé \( r<1\).

	Les points \( 0\) et \( 1\) sont par contre sur la frontière de \( \mathopen[ 0 , 1 \mathclose]\) lorsque cet ensemble est vu comme partie de l'espace métrique \( \eR\).
\end{example}

\begin{lemma}[Passage de douane\cite{ooDKEWooFqlDyN,ooWBUCooAdPjMK}]        \label{LEMooLKWEooItGnkP}
	Dans un espace topologique, toute partie connexe qui rencontre à la fois une partie \( A\) et son complémentaire rencontre nécessairement la frontière de \( A\).
\end{lemma}

\begin{proof}
	Nommons \( C\) la partie connexe qui intersecte \( A\) et \( A^c\). Les ouverts \( \Int(A)\) et \( X\setminus \bar A\) ne peuvent pas recouvrir \( C\) parce que ce sont deux ouverts disjoints alors que \( C\) est connexe (voir la définition~\ref{DefIRKNooJJlmiD} de la connexité). Donc \( C\) doit contenir des points qui sont dans \( \bar A\) mais pas dans \( \Int(A)\). C'est-à-dire des points de \( \partial A\).
\end{proof}

On vérifiera que les notations et les dénominations sont cohérentes en prouvant la proposition suivante.
\begin{proposition}Pour \( \epsilon > 0\),
	\begin{enumerate}
		\item l'adhérence de \( B(x,\epsilon)\) est \( \bar B(x,\epsilon)\),
		\item l'intérieur de \( \bar B(x,\epsilon)\) est \( B(x,\epsilon)\),
		\item la boule ouverte \( B(x,\epsilon)\) est un ouvert,
		\item la boule fermée \( \bar B(x,\epsilon)\) est un fermé.
	\end{enumerate}
\end{proposition}

Nous avons également les liens suivants entre intérieur, adhérence, ouvert, fermé et passage au complémentaire (noté \( {}^c\))~:
\begin{proposition}
	Si \( A \subset \eR^n\) et \( A^c = \eR^n\setminus A\), nous
	avons
	\begin{enumerate}
		\item \( (\Int A)^c = \Adh (A^c)\) et \( (\Adh A)^c = \Int
		      (A^c)\),
		\item \( A\) est ouvert si et seulement si \( A^c\) est fermé,
		\item \( \Int A\) est le plus grand ouvert contenu dans \( A\),
		\item \( \Adh A\) est le plus petit fermé contenant \( A\),
	\end{enumerate}
\end{proposition}

\begin{example} \label{ExBFLooUNyvbw}
	Il n'est en général pas vrai que \( \overline{ A\cap B }=\bar A\cap \bar B\). Par exemple si \( A=\mathopen[ 0 , 1 [\) et \( B=\mathopen] 1 , 2 \mathclose]\). Dans ce cas, \( A\cap B=\emptyset\) alors que \( \bar A\cap\bar B=\{ 1 \}\).
\end{example}

%---------------------------------------------------------------------------------------------------------------------------
\subsection{Boules ouvertes, fermées, sphères}
%---------------------------------------------------------------------------------------------------------------------------

\begin{definition}      \label{DEFooPDSJooFcUqKH}
	Soit un espace métrique\footnote{Définition \ref{DefMVNVFsX}.} \( (E,d)\).
	\begin{enumerate}
		\item
		      Nous nommons \defe{boule fermée}{boule fermée} la fermeture de la boule ouverte, c'est-à-dire les parties de la forme \( \overline{ B(a,r) }\).
		\item
		      La \defe{sphère}{sphère} de centre \( a\in E\) et de rayon \( r\in \eR^+\) est la frontière\footnote{Frontière, définition \ref{DEFooACVLooRwehTl}.} de la boule : \( S(a,r)=\partial B(a,r)\).
	\end{enumerate}
\end{definition}

\begin{normaltext}
	Les différences entre boules ouverts, fermées et sphères sont très importantes. D'abord, les \emph{boules} sont pleines tandis que la \emph{sphère} est creuse. En comparant à une pomme, la boule ouverte serait la pomme «sans la peau», la boule fermée serait «avec la peau» tandis que la sphère serait seulement la peau.
\end{normaltext}


%---------------------------------------------------------------------------------------------------------------------------
\subsection{Espace localement compact}
%---------------------------------------------------------------------------------------------------------------------------

\begin{definition}  \label{DefEIBYooAWoESf}
	Un espace topologique est \defe{localement compact}{compact!localement} si tout élément possède un voisinage compact.
\end{definition}

\begin{lemma}       \label{LEMooAXESooYvyesg}
	Si \( X\) est un espace topologique localement compact et si \( K\) est compact dans \( X\), il existe un ouvert \( V\) tel que \( K\subset V\) et \( \bar V\) est compact.
\end{lemma}

\begin{proof}
	Pour chaque \( x\in K\) nous considérons un voisinage compact \( L_x\). Nous considérons un ouvert \( A_x\) dans \( L_x\) contenant \( x\). Vu que les \( A_x\) recouvrent \( K\), nous avons un sous-recouvrement fini, c'est à dire un choix de \( \{ x_1,\ldots,x_n \}\) tels que
	\begin{equation}
		K\subset\bigcup_{i=1}^nA_{x_i}.
	\end{equation}
	En posant \( V=\bigcup_{i=1}^nA_{x_i}\) nous avons un ouvert qui vérifie \( K\subset V\).

	Montrons que \( \bar V\) est compacte.

	En tant qu'union finie de compacts, la partie \( L=\bigcup_{i=1}^{n}L_{x_i}\) est compacte. Étant donné que \( X\) est séparé, \( L\) est fermé (lemme \ref{LemnAeACf}\ref{ITEMooAZWVooLyPDeY}). Par ailleurs \( \bar V\) est l'intersection de tous les fermés contenant \( V\) (lemme \ref{LEMooILNCooOFZaTe}\ref{ITEMooAZWVooLyPDeY}). Donc \( \bar V\subset L\). La partie \( \bar V\) est donc un fermé dans un compact. Elle est donc compacte par le lemme \ref{LemnAeACf}\ref{ITEMooNKAKooQoNddr}.
\end{proof}


\begin{proposition}[\cite{BIBooMistral,MonCerveau}]	\label{PROPooACXGooHcQpxF}
	Soient un espace topologique \( X\) séparé et localement compact, un ouvert \( U\) et \( x\in U\). Il existe un voisinage ouvert \( V\) de \( x\) tel que
	\begin{enumerate}
		\item
		      \( \bar V\) soit compact,
		\item
		      \( \bar V\subset U\).
	\end{enumerate}
\end{proposition}

\begin{proof}
	Vu que \( X\) est localement compact, nous considérons un voisinage compact \( K\) de \( x\). Plus précisément, nous considérons un ouvert \( A\) et un compact \( K\) tels que
	\begin{equation}
		x\in A\subset K.
	\end{equation}

	Vu que \( U\) est ouvert, \( X\setminus U\) est fermé et donc
	\begin{equation}
		K\setminus U=K\cap (X\setminus U)
	\end{equation}
	est fermé dans le compact \( K\) et est donc compact\footnote{Nous avons utilisé \ref{LemnAeACf} deux fois : une fois pour dire que \( K\) est fermé et une fois pour dire qu'un fermé dans un compact est compact.}.

	Deux possibilités : \( K\setminus U\) est vide ou non.
	\begin{subproof}
		\spitem[Si \( K\setminus U=\emptyset\)]
		%-----------------------------------------------------------
		Dans ce cas \( K\subset U\). En posant \( V=A\) on est bon parce que
		\begin{equation}
			\bar V=\bar A\subset K\subset U.
		\end{equation}
		Notez que \( \bar A\) est l'intersection de tous les fermés contenant \( A\). C'est pour ça que \( \bar A\subset K\).

		De plus \( \bar V\) est compact en tant que fermé dans le compact \( K\).

		\spitem[Si \( K\setminus U\neq \emptyset\)]
		%-----------------------------------------------------------
		Donc \( K\setminus U\) est un compact qui ne contient pas \( x\). Vu que les ouverts séparent les points et les compacts (lemme \ref{LEMooNIJEooMdvkuA}), nous avons des ouverts disjoints \( V_0\) autour de \( x\) et \( W\) autour de \( K\setminus U\) tels que \( V_0\cap W=\emptyset\). Nous posons
		\begin{equation}
			V=V_0\cap A.
		\end{equation}
		C'est encore un ouvert comme intersection d'ouverts. Nous avons
		\begin{equation}
			V\subset A\subset K.
		\end{equation}
		Prouvons que \( \bar V\subset U\). Pour cela nous supposons par l'absurde que \( y\in \bar V\) avec \( y\not\in U\). Donc
		\begin{equation}
			y\in K\setminus U\subset W
		\end{equation}
		Mais \( V\) et \( W\) sont des ouverts d'intersection vide, donc \( \bar V\cap W\) est également vide (proposition \ref{PROPooMDQDooBYjvXO}). Donc \( y\in W\) n'est pas possible. Contradiction. Nous en déduisons que \( \bar V\subset U\).

		Encore une fois \( \bar V\) est compact parce que \( V\subset K\).
	\end{subproof}
\end{proof}

\begin{proposition}[\cite{MonCerveau}]	\label{PROPooAXUEooQxRGmr}
	Soient des ouverts disjoints \( U\) et \( V\) d'un espace topologique. Alors
	\begin{subequations}
		\begin{align}
			U\cap\overline{ V} & = \emptyset \\
			\bar U\cap V       & = \emptyset
		\end{align}
	\end{subequations}
\end{proposition}

\begin{proof}
	Soit \( x\in U\). Étant donné que \( U\) est ouvert, nous avons un voisinage ouvert \( A\) de \( x\) tel que \( A\subset U\). Un tel voisinage n'intersecte donc pas \( V\). Donc \( x\not\in \bar V\).
\end{proof}

\begin{lemma}[\cite{MonCerveau,BIBooMistral}]       \label{LEMooKYMKooPxZjWN}
	Soient un espace topologique séparé localement compact \( X\), un compact \( K\) et un ouvert \( U\) tel que \( K\subset U\). Il existe un ouvert \( V\) tel que
	\begin{enumerate}
		\item
		      \( \bar V\) est compact,
		\item
		      \( K\subset V\subset\bar V\subset U\).
	\end{enumerate}
\end{lemma}

\begin{proof}
	Pour chaque \( x\in K\) la proposition \ref{PROPooACXGooHcQpxF} nous permet de considérer un ouvert \( V_x\) tel que \( \bar V_x\) est compact et \( \bar V_x\subset U\). Compactité, sous-recouvrement fini, tout ça, bref on a \( x_1,\ldots,x_n\in K\) tels que en posant \( V_i=V_{x_i}\), et
	\begin{equation}
		V=\bigcup_{i=1}^nV_i,
	\end{equation}
	on ait \( K\subset V\). La proposition \ref{PROPooEQYEooZHXKnL} nous dit que
	\begin{equation}
		\bar V=\bigcup_{i=1}^n\bar V_i.
	\end{equation}
	En tant qu'union finie de compacts, \( \bar V\) est compacte par le lemme \ref{LEMooFJZDooSxYWVW}. De plus pour tout \( x\in K\) nous avons \( \bar V_x\subset U\). Donc \( \bar V\subset U\).
\end{proof}


\begin{proposition}[\cite{BIBooECOAooBOWCHL,BIBooRIYYooKAxXfm,BIBooWFFBooUFQlGq}]	\label{PROPooVLOZooLlYNNa}
	Soit un espace topologique séparé \( X\). Les faits suivants sont équivalents :
	\begin{enumerate}
		\item		\label{ITEMooZLNUooHcTakl}
		      \( X\) est localement compact\footnote{Définition \ref{DefEIBYooAWoESf}.}.
		\item		\label{ITEMooYMKMooUkcNzd}
		      Tout éléments de \( X\) a un voisinage compact.
		\item		\label{ITEMooXYYHooTCZMJr}
		      Tout point de \( X\) possède une base de voisinages formés d'ouverts à fermeture compactes.
		\item		\label{ITEMooBTECooDDsTuo}
		      Tout point de \( X\) possède un voisinage ouvert à fermeture compacte.
	\end{enumerate}
\end{proposition}

\begin{proof}
	Nous avons \ref{ITEMooZLNUooHcTakl} \( \Leftrightarrow\) \ref{ITEMooYMKMooUkcNzd} parce que c'est la définition de localement compact. À part ça, on va faire un plusieurs parties.
	\begin{subproof}
		\spitem[\ref{ITEMooYMKMooUkcNzd} \( \Rightarrow\) \ref{ITEMooBTECooDDsTuo}]
		%-----------------------------------------------------------
		Soit \( X\in X\). Par \ref{ITEMooYMKMooUkcNzd} nous avons un ouvert \( A\) et un compact \( K\) tels que \( x\in A\subset K\). La partie \( \bar A\) est un fermé dans le compact \( K\) et est donc compacte par le lemme \ref{LemnAeACf}\ref{ITEMooNKAKooQoNddr}.

		\spitem[\ref{ITEMooBTECooDDsTuo} \( \Rightarrow\) \ref{ITEMooYMKMooUkcNzd}]
		%-----------------------------------------------------------
		Soit \( x\in X\). Si \( A\) est le voisinage ouvert à fermeture compacte, alors \( \bar A\) est un voisinage compact.

		\spitem[\ref{ITEMooYMKMooUkcNzd} \( \Rightarrow\) \ref{ITEMooXYYHooTCZMJr}]
		%-----------------------------------------------------------
		Soit un ouvert \( U\) contenant \( x\). Nous allons montrer qu'il existe un ouvert \( A\) autour de \( x\) tel que \( \bar A\) est compact et \( A\subset U\). Soit un point \( x\in X\). Par \ref{ITEMooYMKMooUkcNzd}, il possède un voisinage compact \( K\).

		Nous considérons la frontière \( \partial(U\cap K)\). C'est une partie fermée de \( X\) par la proposition \ref{PROPooTHOPooCOGmZD}\ref{ITEMooCSHKooHibSre}. Et comme \( \partial(U\cap K)\) est contenu dans \( K\), il est un fermé dans un compact. Donc \( \partial(U\cap K)\) est compact.

		Si \( y\in\partial(U\cap K)\), alors \( y\neq x\) parce que \( U\) est un ouvert et n'intersecte pas sa frontière\footnote{Proposition \ref{PROPooFHURooGXcekV}.}. Étant donné que \( X\) est séparé, pour chaque \( y\in \partial(U\cap K)\), nous trouvons des voisinages ouverts \( V_y\) de \( y\) et \( W_y\) de \( x\) tels que \( V_y\cap W_y=\emptyset\).

		Notons que \( K\cap U\) contient un ouvert autour de \( x\) parce que \( K\) en contient un et que \( U\) en est un. Nous pouvons donc choisir \( W_y\subset K\cap U\). L'ensemble de parties \( \{ V_y \}_{y\in\partial(U\cap K)}\) est un recouvrement du compact \( \partial(U\cap K)\) par des ouverts. Nous pouvons donc choisir \( y_1,\ldots,y_n\in\partial(U\cap K)\) tel que
		\begin{equation}
			\partial(U\cap K)\subset \bigcup_{k=1}^nV_{y_k}.
		\end{equation}
		Donc \(A= W_{y_1}\cap\ldots\cap W_{y_n}\)  est un voisinage de \( x\) qui :
		\begin{enumerate}
			\item
			      est inclut dans \( U\cap K\),
			\item
			      n'intersecte aucun des \( V_{y_i}\),
			\item
			      n'intersecte donc pas \( \partial(U\cap K)\).
		\end{enumerate}
		La proposition \ref{PROPooYHIGooMDVKNy} dit que la fermeture \(\bar A= \overline{W_{y_1}\cap\ldots \cap W_{y_n}}\) est encore dans \( U\cap K\). Et comme \( U\cap K\) est compact, nous déduisons que \( \bar A\) est compact.

		\spitem[\ref{ITEMooXYYHooTCZMJr} \( \Rightarrow\) \ref{ITEMooBTECooDDsTuo}]
		%-----------------------------------------------------------
		Si on a une base de voisinages, à fortiori, on a un voisinage.
	\end{subproof}
\end{proof}



%---------------------------------------------------------------------------------------------------------------------------
\subsection{Compactifié d'Alexandrov}
%---------------------------------------------------------------------------------------------------------------------------

\begin{propositionDef}[\cite{ooEDBNooKshWkw}]       \label{PROPooHNOZooPSzKIN}
	Soit un espace topologique séparé localement compact\footnote{Définition \ref{DefEIBYooAWoESf}.} \( X\). Nous considérons un élément \( \omega\notin X\) et l'ensemble \( \hat X =\ X\cup\{ \omega \}\). Nous nommons «ouverts de \( \hat X\)» les parties suivantes :
	\begin{itemize}
		\item les ouverts de \( X\),
		\item les parties de la forme \( A\cup\{ \omega \}\) avec \( X\setminus A\) compact dans \( X\).
	\end{itemize}
	Alors \( \hat X\) est un espace topologique compact (cela justifie le nom «ouvert» donné aux parties sus-définies).
\end{propositionDef}

\begin{proof}
	La première chose à faire est de prouver que \( \hat X\) est bien un espace topologique (définition \ref{DefTopologieGene}). Nous notons \( \tau\) la topologie sur \( X\) et \( \hat\tau\) l'ensemble des «ouverts» de \( \hat X\). Le but est de prouver que \( \hat \tau\) est une topologie.
	\begin{subproof}
		\spitem[L'espace lui-même]
		\( \hat X\in \hat\tau\) parce que \( \hat X=X\cup \{ \omega \}\) et que \( X\setminus X=\emptyset\) est compact.
		\spitem[Le vide]
		\( \emptyset\in \tau\subset \hat \tau\).
		\spitem[Union quelconque]
		Soient \( A_i\) (\( i\in I\)) des éléments de \( \hat\tau\). Nous posons \( I_1=\{ i\in I\tq A_i\subset X \}\) et \( I_2=I\setminus I_1\). Nous avons
		\begin{equation}
			\bigcup_{i\in I}A_i=\big( \bigcup_{i\in I_1}A_i \big)\cup \big( \bigcup_{i\in I_2}A_i \big)=B\cup\big( \bigcup_{i\in I_2}B_i\cup\{ \omega \} \big)
		\end{equation}
		où \( B\) et les \( B_i\) sont des ouverts de \( X\) tels que \( X\setminus B_i\) est compact dans \( X\). Nous récrivons ça sous la forme
		\begin{equation}
			\bigcup_{i\in I}A_i=B\cup\big( \bigcup_{i\in I_2}B_i \big)\cup\{ \omega \}.
		\end{equation}
		La question est de savoir si
		\begin{equation}
			X\setminus\Big( B\cup\big( \bigcup_{i\in I_2}B_i \big) \Big)
		\end{equation}
		est compact dans \( X\). Un peu de réécriture :
		\begin{equation}
			X\setminus\Big( B\cup\big( \bigcup_{i\in I_2}B_i \big) \Big)=(X\setminus B)\cap X\setminus\big( \bigcup_{i\in I_2}B_i \big)=(X\setminus B)\cap\big( \bigcap_{i\in I_2}(X\setminus B_i) \big).
		\end{equation}
		La partie \( X\setminus B\) est fermée dans \( X\) parce que \( B\) est ouverte. La proposition \ref{PROPooQWHSooXeJOkT} dit qu'une intersection de compacts est compacte (parce que \( X\) est séparé). Nous sommes donc en présence de l'intersection entre un compact et un fermé.

		Tout compact d'un espace séparé est fermé\footnote{Lemme \ref{LemnAeACf}\ref{ITEMooAZWVooLyPDeY}.}. Donc nous sommes en présence de l'intersection de deux fermés. Donc
		\( (X\setminus B)\cap\big( \bigcap_{i\in I_2}(X\setminus B_i) \big)\) est fermé. Mais c'est contenu dans le compact \( \bigcap_{i\in I_2}(X\setminus B_i)\). Fermé dans un compact, donc compact (lemme \ref{LemnAeACf}).
		\spitem[Intersection finie]
		Nous considérons les «ouverts» \( (A_i)_{i=1,\ldots, n}\) de \( \hat X\). Si ce sont tous des ouverts de \( X\), l'intersection est un ouvert de \( X\) et on est bon.

		Supposons que tous les \( A_i\) soient de la forme \( A_i=B_i\cup\{ \omega \}\) avec \( X\setminus B_i\) compact. Alors
		\begin{equation}
			\bigcap_{i=1}^n (B_i\cup\{ \omega \})=\left( \bigcap_{i=1}^nB_i \right)\cup\{ \omega \}
		\end{equation}


		Mais le lemme \ref{LEMooHRKAooRskzQL} (appliqué un nombre fini de fois) donne
		\begin{equation}
			X\setminus\left( \bigcap_{i=1}^nB_i \right)=\bigcup_{i=1}^n(X\setminus B_i)
		\end{equation}
		qui est compact en tant qu'union finie de compacts\footnote{Lemme \ref{LEMooFJZDooSxYWVW}.}.

		Enfin, nous supposons que les \( A_i\) sont un mélange des deux types, nous les séparons entre ceux qui sont directement des ouverts de \( X\) et les autres :
		\begin{equation}
			A_i=\begin{cases}
				B_i                 & \text{si } i\leq q  \\
				C_i\cup\{ \omega \} & \text{si }q<i\leq n
			\end{cases}
		\end{equation}
		où \( B_i\) sont ouverts et \( C_i\) sont des parties de \( X\) telles que \( X\setminus C_i\) est compacte.

		Nous avons
		\begin{subequations}
			\begin{align}
				\bigcap_{i=1}^nA_i & =\left( \bigcap_{i=1}^qB_i \right)\cap\left( \bigcap_{i=q+1}^n(C_i\cup\{ \omega \}) \right) \\
				                   & =B\cap\left( \bigcap_{i=q+1}^nC_i \right).
			\end{align}
		\end{subequations}
		Justifications.
		\begin{itemize}
			\item Nous avons posé \( B\) est l'intersection des \( B_i\).
			\item
			      Vu que \( \omega\) n'est pas dans \( B\), nous pouvons l'oublier dans les \( C_i\cup\{ \omega \}\).
		\end{itemize}
		C'est le moment d'étudier \( E=\bigcap_{i=q+1}^nC_i\). Nous avons
		\begin{equation}
			X\setminus E=X\setminus\left( \bigcap_{i=q+1}^nC_i \right)=\bigcup_{i=q+1}^n(X\setminus C_i).
		\end{equation}
		Vu que \( X\setminus C_i\) est compact, il est fermé\footnote{Par le lemme \ref{LemnAeACf}\ref{ITEMooAZWVooLyPDeY} et le fait que nous considérons un espace séparé.}. La partie \( X\setminus E\) est donc fermée comme union finie de fermés\footnote{Par le lemme \ref{LemQYUJwPC}\ref{ItemKJYVooMBmMbG}.}. Et donc \( E\) est ouvert. Et finalement
		\begin{equation}
			\bigcap_{i=1}^nA_i=B\cap E
		\end{equation}
		est un ouvert de \( X\) comme intersection d'ouverts. C'est donc aussi un ouvert de \( \hat X\).

	\end{subproof}
	Nous avons fini de prouver que \( (\hat X, \hat \tau)\) est un espace topologique. Nous montrons à présent que \( \hat X\) est compact.

	Soit \( \{ A_i \}_{i\in I}\) un recouvrement de \( \hat X\) par des ouverts. Pour au moins un \( i_0\in I\) nous avons \( \omega\in A_{i_0}\). Nous posons \( A_{i_0}=B\cup\{ \omega \}\) avec \( X\setminus B\) compact dans \( X\).

	Les ouverts \( \{ A_i \}_{i\in I}\) forment un recouvrement de \( X\setminus B\) par des ouverts. Nous pouvons en extraire un sous-recouvrement fini :
	\begin{equation}
		X\setminus B\subset \bigcup_{i\in I_1}A_i.
	\end{equation}
	Nous avons alors
	\begin{equation}        \label{EQooJQFMooKPXLkh}
		\hat X\subset \bigcup_{i\in I_1\cup\{ i_0 \}}A_i.
	\end{equation}
	Et voilà que \( \hat X\) est recouvert par un nombre fini des \( A_i\). Notez que \eqref{EQooJQFMooKPXLkh} est une égalité, mais nous n'en avons pas besoin.
\end{proof}

\begin{normaltext}
	Oh bien entendu, les plus férus de questions embarrassantes demanderont, si \( X\) est l'espace considéré, où prendre ce \( \omega\) ? Quel «objet» existe en-dehors de \( X\) ? Qui m'assure que \( X\) n'est pas tellement grand que tout est dedans ? Le fait est qu'il n'existe pas d'ensemble contenant tous les ensembles (c'est le corolaire \ref{CORooZMAOooPfJosM}). Nous pouvons donc toujours trouver un ensemble \( \omega\) qui n'est pas dans \( X\).
\end{normaltext}

En ce qui concerne \( \eR\) auquel nous pouvons attacher deux infinis (\( +\infty\) et \( -\infty\)), ce sera la définition \ref{DEFooRUyiBSUooALDDOa}.

Pour \( \eC\), nous donnerons une caractérisation de la limite en \( \infty\) dans le lemme \ref{LEMooERABooQjLBzW}.




%---------------------------------------------------------------------------------------------------------------------------
\subsection{Continuité séquentielle}
%---------------------------------------------------------------------------------------------------------------------------

\begin{corollary}[Caractérisation séquentielle de la continuité en un point\cite{MonCerveau}]		\label{PropFnContParSuite}
	Une application entre deux espaces topologiques continue en un point y est séquentiellement\footnote{Définition \ref{DefENioICV}.} continue.
\end{corollary}

\begin{proof}
	Soit une application \( f\colon X\to Y\) entre les espaces topologiques \( X\) et \( Y\). Nous supposons que \( f\) est continue en \( a\in X\). Soit une suite convergente \( x_k\stackrel{X}{\longrightarrow}a\). Nous devons prouver que \( f(x_k)\to f(a)\).

	Soit un voisinage \( V\) de \( f(a)\) dans \( Y\). Le fait que \( f\) soit continue en \( a\) signifie\footnote{C'est la définition \ref{DefOLNtrxB} de la continuité en un point.} que \( f(a)\) est une limite de \( f\) en \( a\), c'est-à-dire\footnote{Définition \ref{DefYNVoWBx} d'une limite.} qu'il existe un voisinage \( W\) de \( a\) tel que \( f(W\setminus\{ a \})\subset V\).

	Puisque \( x_k\to a\), il existe \( N\) tel que \( x_k\in W\) pour tout \( k\geq N\). Pour ces valeurs de \( k\), nous avons \( f(x_k)\in V\).

	Nous avons prouvé que pour tout voisinage \( V\) de \( f(a)\) dans \( Y\), il existe \( N\) tel que \( f(x_k)\in V\) dès que \( k\geq N\). Cela signifie exactement que \( f(x_k)\to f(a)\).
\end{proof}

\subsubsection{Les boules, une base de topologie}
%////////////////////////////

\begin{proposition} \label{PropNBSooraAFr}
	Un espace métrique séparable\footnote{Qui possède une partie dense dénombrable, définition~\ref{DefUADooqilFK}.} accepte une base de topologie\footnote{Base de topologie, définition \ref{DEFooLEHPooIlNmpi}.} dénombrable.

	Soit \( A\) dense et dénombrable dans l'espace métrique séparable \( (E,d)\). Si \( \{ a_i \}_{i\in \eN}\) est une énumération de \( A\) et \( \{ r_i \}_{i\in \eN}\) une énumération de \( \eQ\), alors
	\begin{equation}
		\mB=\{ B(a_i,r_j) \}_{i,j\in \eN}
	\end{equation}
	est une base de la topologie\footnote{Définition \ref{DEFooLEHPooIlNmpi}.} de \( E\).
\end{proposition}
\index{base!de topologie!espace métrique}
\index{espace!métrique!base de topologie}
\index{base!de topologie!dénombrable}

\begin{proof}
	Soient \( x\in E\) et \( V\) un voisinage de \( x\). Ce dernier contient une boule \( B(x,r)\) et quitte à prendre \( r\) un peu plus petit nous supposons que \( r\in \eQ\) (existence d'un tel rationnel par le lemme \ref{LemooHLHTooTyCZYL}).

	Soit \( a\in A\) avec \( \| a-x \|<\frac{ r }{ 3 }\) (existe par densité de \( A\) dans \( E\)); nous avons \( B(a,\frac{ 2r }{ 3 })\subset B(x,r)\) parce que si \( y\in B( a,\frac{ 2r }{ 3 } )\) alors
	\begin{equation}
		\| y-x \|\leq \| y-a \|+\| a-x \|<\frac{ 2 }{ 3 }r+\frac{ 1 }{ 3 }r=r.
	\end{equation}
	La seconde inégalité est stricte parce que les boules sont ouvertes. Le tout montre que \( y\in B(x,r)\). Par ailleurs \( x\in B(a,\frac{ 2 }{ 3 }r)\) et nous avons trouvé un élément de \( \mB\) contenant \( x\) tout en étant inclus dans \( V\). Cela prouve que \( \mB\) est bien une base de la topologie de \( E\).
\end{proof}


\begin{remark}      \label{RemIPVLooHUXyeW}
	Il est vite vu que les cubes ouverts forment aussi une base de la topologie de \( \eR^n\). Cela est à mettre en rapport avec le fait que toutes les normes sont équivalentes sur \( \eR^n\) (proposition~\ref{ThoNormesEquiv}).

	% position 13268

	Voir aussi le corolaire~\ref{CorTHDQooWMSbJe} qui donnera tout ouvert comme union de pavés presque disjoints.
\end{remark}

\begin{definition}\label{DefEnsembleBorne}
	Soit \( (X, d) \) un espace métrique. Un sous-ensemble \( A \subset X\) est \defe{borné}{borné} si il existe une boule de \( X\) contenant \( A\).
\end{definition}

\begin{proposition}     \label{PROPooJIOAooWqzKMu}
	Toute réunion finie de parties bornés est bornée. Toute partie d'un ensemble borné est un ensemble borné.
\end{proposition}

\begin{proof}
	Soient des parties bornées \( \{ A_i \}_{i=1,\ldots,n}\). Si \( m_i\) est un majorant de \( A_i\), alors \( M=\max\{ m_i \}_{i=1,\ldots,n}\) est un majorant de \( \bigcup_{i=1}^nA_i\). Notons que \( M\) est bien défini par le lemme \ref{LEMooPCRFooXRGrUr}.

	Pour la seconde affirmation, si \( A\) est borné (notons \( m\) un majorant de \( A\)) et si \( B\subset A\), alors \( m\) est un majorant de \( B\).
\end{proof}


%---------------------------------------------------------------------------------------------------------------------------
\subsection{Continuité et compacité}
%---------------------------------------------------------------------------------------------------------------------------

Un résultat important dans la théorie des fonctions sur les espaces vectoriels normés est qu'une fonction continue sur un compact est bornée et atteint ses bornes. Ce résultat sera énormément utilisé pour trouver des maximums et minimums de fonctions. Le théorème exact est le suivant.

\begin{lemma}[de Lebesgue\cite{AntoniniAndAl-EspacesMetriquesCompacts}]    \label{LemQFXOWyx}
	Soit \( (X,d)\) un espace métrique tel que toute suite ait une sous-suite convergente à l'intérieur de l'espace. Si \( \{ V_i \}\) est un recouvrement par des ouverts de \( X\), alors il existe \( \epsilon\) tel que pour tout \( x\in X\), nous ayons \( B(x,\epsilon)\subset V_i\) pour un certain \( i\).
\end{lemma}

\begin{proof}
	Par l'absurde, nous supposons que pour tout \( n\), il existe un \( x_n\in X\) tel que la boule \( B(x_n,\frac{1}{ n })\) n'est contenue dans aucun des \( V_i\). De ces \( x_n\), nous extrayons une sous-suite convergente (que nous nommons encore \( (x_n)\)) et nous posons \( x_n\to x\). Pour \( n\) assez grand (\( \frac{1}{ n }<\epsilon\)) nous avons \( x_n\in B(x,\epsilon)\), donc tous les \( x_n\) suivants sont dans le \( V_i\) qui contient \( x\).
\end{proof}

\begin{lemma}[\cite{AntoniniAndAl-EspacesMetriquesCompacts}]   \label{LemMGQqgDG}
	Soit \( (X,d)\) un espace métrique tel que toute suite possède une sous-suite convergente. Pour tout \( \epsilon>0\), il existe un ensemble fini \( \{ x_i \}_{i\in I}\) tel que les boules \( B(x_i,\epsilon)\) recouvrent \( X\).
\end{lemma}

\begin{proof}
	Soit par l'absurde un \( \epsilon>0\) contredisant le lemme. Il n'existe pas de parties finies de \( X\) autour des points desquels les boules de taille \( \epsilon\) recouvrent \( X\).

	Nous construisons par récurrence une suite ne possédant pas de sous-suite convergente. Le premier terme, \( x_0\) est pris arbitrairement dans \( X\). Ensuite si nous avons déjà \( N\) termes de la suite, nous savons que les boules de rayon \( \epsilon\) centrées sur les points \( \{ x_i \}_{i=1,\ldots, N}\) ne recouvrent pas \( X\). Donc nous prenons \( x_{N+1}\) hors de l'union de ces boules.

	Ainsi nous avons une suite \( (x_n)\) dont tous les termes sont à distance plus grande que \( \epsilon\) les uns des autres. Une telle suite ne peut pas contenir de sous-suite convergente. Contradiction.
\end{proof}

\begin{theorem}[Bolzano-Weierstrass\cite{AntoniniAndAl-EspacesMetriquesCompacts}, thème \ref{THEMEooQQBHooLcqoKB}]\label{ThoBWFTXAZNH}
	Un espace métrique est compact si et seulement si toute suite admet une sous-suite qui converge à l'intérieur de l'espace.
\end{theorem}
\index{théorème!Bolzano-Weierstrass}
\index{Bolzano-Weierstrass!espaces métriques}
\index{compacité}

\begin{proof}
	Soient \( X\) un espace métrique compact et \( (x_n)\) une suite dans \( X\). Nous considérons la suite de fermés emboîtés
	\begin{equation}
		X_n=\overline{ \{ x_k\tq k>n \} }.
	\end{equation}
	Ce sont des fermés ayant la propriété d'intersection finie non vide, et donc la proposition~\ref{PropXKUMiCj} nous dit qu'ils ont une intersection non vide. Un élément de cette intersection est automatiquement un point d'accumulation de la suite\footnote{Définition \ref{DEFooGHUUooZKTJRi}.}.

	Nous passons à l'autre sens. Nous supposons que toute suite dans \( X\) contient une sous-suite convergente, et nous considérons \( \{ V_i \}_{i\in I}\), un recouvrement de \( X\) par des ouverts. Par le lemme~\ref{LemQFXOWyx}, nous considérons un \( \epsilon\) tel que pour tout \( x\), il existe un \( i\in I\) avec \( B(x,\epsilon)\subset V_i\). Par le lemme~\ref{LemMGQqgDG}, nous considérons un ensemble fini \( \{ y_i \}_{i\in A}\) tel que les boules \( B(y_i,\epsilon)\) recouvrent \( X\).

	Par construction, chacune de ces boules \( B(y_i,\epsilon)\) est contenue dans un des ouverts \( V_i\). Nous sélectionnons donc parmi les \( V_i\) le nombre fini qu'il faut pour recouvrir les \( B(y_i,\epsilon)\) et donc pour recouvrir \( X\).
\end{proof}

\begin{example}[Non compacité de la boule unité en dimension infinie]\label{ExEFYooTILPDk}
	Le théorème de Bolzano-Weierstrass permet de voir tout de suite que la boule unité n'est pas compacte dans un espace vectoriel de dimension infinie : la suite des vecteurs de base ne possède pas de sous-suite convergente.
\end{example}


Le théorème de Bolzano–Weierstrass~\ref{ThoBWFTXAZNH} a l'importante conséquence suivante.
\begin{theorem}[Weierstrass]		\label{ThoWeirstrassRn}
	Une fonction continue à valeurs réelles définie sur un compact est bornée et atteint ses bornes.
\end{theorem}
\index{théorème!Weierstrass}
\index{compact!et fonction continue}

\begin{proof}
	Soient \( K\) un compact et \( f\colon K\to \eR\) une fonction continue. Nous désignons par \( A\) l'ensemble des valeurs prises par \( f\) sur \( K\) :
	\begin{equation}
		A=f(K)=\{ f(x)\tq x\in K \}.
	\end{equation}
	Nous considérons le supremum \( M=\sup A=\sup_{x\in K}f(x)\) avec la convention suivante : si \( A\) n'est pas borné supérieurement, nous posons \( M=\infty\) (voir définition~\ref{DefSupeA}).

	Nous allons maintenant construire une suite \( (x_n)\) de deux façons différentes selon que \( M=\infty\) ou non.
	\begin{enumerate}
		\item
		      Si \( M=\infty\), nous choisissons, pour chaque \( n\in\eN\), un \( x_n\in K\) tel que \( f(x_n)>n\). C'est certainement possible parce que si \( A\) n'est pas borné, nous pouvons y trouver des nombres aussi grands que nous voulons.
		\item
		      Si \( M\neq\infty\), nous savons que pour tout \( \varepsilon\), il existe un \( y\in A\) tel que \( y>M-\varepsilon\). Pour chaque \( n\), nous choisissons donc \( x_n\in K\) tel que \( f(x_n)>M-\frac{1}{ n }\).
	\end{enumerate}
	Quel que soit le cas dans lequel nous sommes, la suite \( (x_n)\) est une suite dans \( K\) qui est compact, et donc nous pouvons en extraire une sous-suite convergente à l'intérieur de \( K\) par le théorème de Bolzano-Weierstrass~\ref{ThoBWFTXAZNH}. Afin d'alléger la notation, nous allons noter \( (x_n)\) la sous-suite convergente. Nous avons donc
	\begin{equation}
		x_n\to x\in K.
	\end{equation}
	Par la proposition~\ref{PropFnContParSuite}, nous savons que \( f\) prend en \( x\) la valeur
	\begin{equation}
		f(x)=\lim_{n\to \infty} f(x_n).
	\end{equation}
	Donc \( f(x)<\infty\). Évidemment, si nous avions été dans le cas où \( M=\infty\), la suite \( x_n\) aurait été choisie pour avoir \( f(x_n)>n\) et donc il n'aurait pas été possible d'avoir \( \lim_{n\to \infty} f(x_n)<\infty\). Nous en concluons que \( M<\infty\), et donc que \( f\) est bornée sur \( K\).

	Afin de prouver que \( f\) atteint sa borne, c'est-à-dire que \( M\in A\), nous considérons les inégalités
	\begin{equation}
		M-\frac{1}{ n }<f(x_n)\leq M.
	\end{equation}
	En passant à la limite \( n\to \infty\), ces inégalités deviennent
	\begin{equation}
		M\leq f(x)\leq M,
	\end{equation}
	et donc \( f(x)=M\), ce qui prouve que \( f\) atteint sa borne \( M\) au point \( x\in K\).
\end{proof}

\begin{lemma}[\cite{MonCerveau}]       \label{LEMooQLVAooICaPvR}
	Soient des compacts \( A,B\) et une fonction continue \( f\colon A\times B\to \eR\). Alors
	\begin{equation}
		\sup_{(x,y)\in A\times B}| f(x,y) |=\sup_{x\in A}\big( \sup_{y\in B}| f(x,y) | \big).
	\end{equation}
\end{lemma}

\begin{proof}
	Pour chaque \( x\in A \), la fonction \( f_x\colon B\to \eR\) donnée par \( f_x(y)=| f(x,y) |\) est continue et atteint donc sa borne\footnote{Théorème \ref{ThoWeirstrassRn}.} en \( y_M(x)\). Notons que cela ne définit pas de façon univoque \( y_M(x)\) parce que \( f_x\) peut atteindre son maximum en plusieurs points. L'important est que pour tout \( x\), le nombre \( | f\big( x,y_M(x) \big) |\) ne dépend pas du choix de \( y_M(x)\) parmi les \( y\) qui réalisent le maximum.


	Notons \( (x_0,y_0)\) un point de \( A\times B\) sur lequel \( | f |\) réalise son maximum\footnote{Encore une fois, ce point n'est pas déterminé de façon unique par cette propriété.} :
	\begin{equation}        \label{EQooDDXDooVsnlKG}
		\sup_{(x,y)\in A\times B}| f(x,y) |=| f(x_0,y_0) |.
	\end{equation}


	Nous avons d'une part
	\begin{equation}
		\sup_{x\in A}\big( \sup_{y\in B}| f(x,y) | \big)=\sup_{x\in A}| f\big( x,y_M(x) \big) |\leq | f(x_0,y_0) |
	\end{equation}

	Et d'autre part, quelques calculs avec justifications en-dessous :
	\begin{subequations}        \label{SUBEQooPYJPooBJyEgN}
		\begin{align}
			\sup_{x\in A}\big( \sup_{y\in B}| f(x,y) | \big) & \leq \sup_{x\in A}\sup_{y\in B}| f(x_0,y_0) |  \label{SUBEQooAKPOooPdkvMJ} \\
			                                                 & =| f(x_0,y_0) |                                                            \\
			                                                 & \leq | f\big(x_0,y_M(x_0)\big) |       \label{SUBEQooDYVPooUgOpfD}         \\
			                                                 & \leq \sup_{x\in A}| f\big(x,y_M(x)\big) |  \label{SUBEQooVOFAooNtzSpt}     \\
			                                                 & \leq \sup_{x\in A}\big( \sup_{y\in B}| f(x,y) | \big).
		\end{align}
	\end{subequations}
	Justifications.
	\begin{itemize}
		\item   Pour \eqref{SUBEQooAKPOooPdkvMJ}. Le point \( (x_0,y_0)\) est un maximum de \( | f |\).
		\item Pour \eqref{SUBEQooDYVPooUgOpfD}. \( y_M\) est définie pour maximiser, en fonction de \( x\), la quantité \( | f(x, y_M(x)) |\).
		\item Pour \eqref{SUBEQooVOFAooNtzSpt}. Au lieu de conserver la valeur \( x_0\) fixé, nous prenons le maximum sur tous les \( x\) possibles.
	\end{itemize}
	Vu que les premiers et derniers termes des inégalités \eqref{SUBEQooPYJPooBJyEgN} sont égaux, toutes les inégalités sont en réalité des égalités. En particulier, en reprenant \eqref{EQooDDXDooVsnlKG},
	\begin{equation}
		\sup_{(x,y)\in A\times B}| f(x,y) |=| f(x_0,y_0) |=\sup_{x\in A}\big( \sup_{y\in B}| f(x,y) | \big).
	\end{equation}
\end{proof}

%---------------------------------------------------------------------------------------------------------------------------
\subsection{Distance à un ensemble}
%---------------------------------------------------------------------------------------------------------------------------

\begin{definition}      \label{DEFooGNNUooFUZINs}
	Si \( A\) est une partie de l'espace métrique \( (X,d)\), et si \( b\in X\), nous définissons
	\begin{equation}
		d(b,A)=\inf_{y\in A}d(b,y).
	\end{equation}
\end{definition}


\begin{lemma}[\cite{MonCerveau}]        \label{LEMooEQIZooLpsbOe}
	Soit un fermé \( F\) de l'espace métrique \( (X,d)\). Si \( a\in X\) vérifie \( d(a,F)=0\), alors \( a\in F\).
\end{lemma}

\begin{proof}
	Supposons que \( d(a,F)=0\), c'est-à-dire que \( \inf_{x\in F} d(a,x) =0\). Il existe donc une suite \( (x_k)\) dans \( F\) telle que \( d(a,x_k)\to 0\).

	Cela signifie que \( x_k\stackrel{(X,d)}{\longrightarrow}a\). La proposition \ref{PROPooBBNSooCjrtRb} nous dit alors que \( a\in F\).
\end{proof}


\begin{example}[Pas avec un ouvert]
	En prenant l'ouvert \( A=\mathopen] 0 , 1 \mathclose[\) dans \( \eR\) nous avons \( d(0,A)=0\), alors que \( 0\) n'est pas dans \( A\).
\end{example}

\begin{lemma}[\cite{MonCerveau}]    \label{LEMooJNRTooZyKiFC}
	Soient un espace métrique \( (X,d)\) ainsi qu'une partie \( A\subset X\). Soit \( r>0\). La partie
	\begin{equation}
		\mO=\{ x\in X\tq d(x,A) < r\}
	\end{equation}
	est ouverte.
\end{lemma}

\begin{proof}
	Soit \( y\in \mO\); nous avons \( d(y,A)<r\). Autrement dit,
	\begin{equation}
		\inf_{a\in A}d(y,a)<r
	\end{equation}
	et donc il existe \( a\in A\) tel que \( d(y,a)<r\). Soit \( \delta=d(y,a)<r\). Nous montrons à présent que \( B(y,r-\delta)\) est dans \( \mO\). En effet si \( z\in B(y,r-\delta)\), alors
	\begin{equation}
		d(z,a)\leq d(z,y)+d(y,a)<r-\delta+\delta=r.
	\end{equation}
\end{proof}


\begin{proposition}[Inégalité triangulaire point-partie\cite{BIBooTHPCooHuDKOy}]	\label{PROPooTXYMooDGnfKB}
	Soit un espace métrique\footnote{Espace métrique, définition \ref{DefMVNVFsX}.} \( (X,d)\) et soit une partie non vide \( A\subset X\). Pour tout \( x,y\in A\) nous avons :
	\begin{equation}
		d(x,A)\leq d(x,A)+d(x,y)
	\end{equation}
\end{proposition}

\begin{proof}
	Par l'inégalité triangulaire pour tout \( a\in A\) nous avons \( d(x,a)\leq d(x,y)+d(y,a)\) et en arrangeant :
	\begin{equation}
		d(y,a)\geq d(x,a)-d(x,y).
	\end{equation}
	Et maintenant on peut calculer deux secondes :
	\begin{subequations}
		\begin{align}
			d(y,A) & =\sup_{a\in A}d(y,a)                        \\
			       & \geq \sup_{a\in A}\big( d(x,a)- d(x,y)\big) \\
			       & =\big( \sup_{a\in A}d(x,a) \big)-d(x,y)     \\
			       & =d(x,A)-d(x,y).
		\end{align}
	\end{subequations}
\end{proof}


\begin{lemma}[\cite{MonCerveau,BIBooXQVFooDewppl}]        \label{LEMooCFGTooIfdcfk}
	Soit une partie \( A\) de l'espace métrique \( (X,d)\). Alors
	\begin{enumerate}
		\item		\label{ITEMooDXSHooPhCXlR}
		      Pour tout \( x,y\in X\) nous avons
		      \begin{equation}
			      | d(x,A)-d(y,A) |\leq d(x,y).
		      \end{equation}
	\end{enumerate}

	\item		\label{ITEMooGNSEooDVZrFp}
	La fonction
	\begin{equation}
		\begin{aligned}
			f\colon X & \to \mathopen[ 0 , \infty \mathclose[ \\
			x         & \mapsto d(x,A)
		\end{aligned}
	\end{equation}
	est continue.
\end{lemma}

\begin{proof}
	Nous utilisons l'inégalité triangulaire point-partie \ref{PROPooTXYMooDGnfKB}. D'abord \( d(x,A)\leq d(x,y)+d(y,A)\) que nous récrivons
	\begin{equation}		\label{EQooCJKBooXdbzts}
		d(x,A)-d(y,A)\leq d(x,y).
	\end{equation}
	Ensuite \( d(y,A)\leq d(x,y)+d(x,A)\) que nous récrivons
	\begin{equation}		\label{EQooKHPCooDSsfTh}
		d(y,A)-d(x,A)\leq d(x,y).
	\end{equation}
	Notez que les équations \eqref{EQooCJKBooXdbzts} et \eqref{EQooKHPCooDSsfTh} sont de la forme
	\begin{subequations}
		\begin{numcases}{}
			u-v\leq \alpha\\
			v-u\leq \alpha,
		\end{numcases}
	\end{subequations}
	dont nous déduisons \( | u-v |\leq \alpha\). Bref, nous avons
	\begin{equation}		\label{EQooNAQCooJzmFoj}
		| d(x,A)-d(y,A) |\leq d(x,y).
	\end{equation}
	Voilà qui prouve déjà \ref{ITEMooDXSHooPhCXlR}.

	Pour prouver \ref{ITEMooGNSEooDVZrFp}, nous considérons \( x_n\stackrel{ X}{\longrightarrow} x\), et écrivons \eqref{EQooNAQCooJzmFoj} pour \( x\) et \( x_n\) : pour tout \( n\) nous avons
	\begin{equation}
		| d(x_n,A)-d(x,A) |\leq d(x_n,x).
	\end{equation}
	Vu que \(d \colon X\times X\to X  \) est continue, en prenant la limite \( n\to \infty\), à droite nous avons zéro\footnote{Parce que \( d\) est continue, proposition \ref{PROPooSZMLooNdtLLj}.} et donc
	\begin{equation}
		\lim_{n\to \infty}| d(x_n,A)-d(x,A) |=0,
	\end{equation}
	ou encore
	\begin{equation}
		\lim_{n\to \infty}| f(x_n)-f(x) |=0,
	\end{equation}
	ce qui signifie que \( f\) est continue en \( x\) parce que la continuité séquentielle implique la continuité par la proposition \ref{PropXIAQSXr}.
\end{proof}



Si \( V\) est une partie d'un espace vectoriel métrique\footnote{Définition \ref{DefMVNVFsX}.} \( E,d\), nous posons
\begin{equation}
	B(V,\delta)=\{ x\in E\tq d(x,V)<\epsilon \}
\end{equation}
où la distance entre \( V\) et \( x\) est définie en \ref{DEFooGNNUooFUZINs}.

\begin{proposition}	\label{PROPooYTMRooExvGrN}
	Soient un ouvert \( U\) de \( \eR^n\) et un compact \( K\subset U\). Alors il existe \( \epsilon>0\) tel que
	\begin{equation}
		K\subset B(K,\epsilon)\subset U.
	\end{equation}
	%TODOooUUYTooQXCVup. Prouver ça.
\end{proposition}


%---------------------------------------------------------------------------------------------------------------------------
\subsection{Espace métrisable}
%---------------------------------------------------------------------------------------------------------------------------

\begin{proposition}[\cite{ooCGEHooVTyTuY}]      \label{PROPooXWBTooCvGLOj}
	Soit un espace topologique métrisable\footnote{Définition \ref{DEFooGLTUooJjaSmI}.} \( X\).
	\begin{enumerate}
		\item   \label{ITEMooOXVRooBsKwuq}
		      Tout fermé de \( X\) est une intersection dénombrable d'ouverts.
		\item
		      Tout ouvert de \( X\) est une union dénombrable de fermés.
	\end{enumerate}
\end{proposition}

\begin{proof}
	Soit une métrique \( d\) compatible avec la topologie de \( X\) et un fermé \( A\). Nous posons
	\begin{equation}
		V_n=\{ x\in X\tq d(x,A)<\frac{1}{ n } \}.
	\end{equation}
	Et juste pour faire simple nous notons \( V_0=X\).
	\begin{subproof}
		\spitem[Les parties \( V_n\) sont ouvertes]
		Soit \( x\in V_n\). Trouvons un voisinage de \( x\) contenu dans \( V_n\) afin de pouvoir encore invoquer le théorème~\ref{ThoPartieOUvpartouv}. D'abord, vu que \( x\in V_n\), il existe \( a\in A\) tel que \( d(x,a)<\frac{ 1 }{ n }\) (ici les inégalités strictes sont importantes).

		Soient \( \epsilon>0\) que nous fixerons plus bas, et \( y\in B(x,\epsilon)\). L'inégalité triangulaire donne
		\begin{equation}
			d(y,a)\leq d(y,x)+d(x,a)<\epsilon+\frac{1}{ n }.
		\end{equation}
		Nous pouvons donc choisir \( \epsilon\) de telle sorte que \( d(y,a)<1/n\). Avec ce \( \epsilon\), nous avons, pour tout \( y\in B(x,\epsilon)\) :
		\begin{equation}
			d(y,A)\leq d(y,a)<\frac{1}{ n }
		\end{equation}
		et donc \( y\in V_n\).
		\spitem[\( A\) est l'intersection des \( V_n\)]
		Nous avons évidemment \( A\subset V_n\) pour tout \( n\). Et d'autre part, si \( a\in\bigcap_{n\in \eN} V_n\) alors \( d(a,A)<\frac{1}{ n }\) pour tout \( n\). Cela implique \( d(a,A)=0\), et donc \( a\in A\) par le lemme \ref{LEMooEQIZooLpsbOe}.
	\end{subproof}

	Ceci démontre le point \ref{ITEMooOXVRooBsKwuq}.

	En ce qui concerne la seconde partie, nous appliquons la première partie au complémentaire. Si \( \mO\) est ouvert, \( \mO^c\) est fermé et
	\begin{equation}
		\mO^c=\bigcap_{n\in \eN}V_n,
	\end{equation}
	ce qui donne immédiatement
	\begin{equation}
		\mO=\bigcup_{n\in \eN}V_n^c
	\end{equation}
	où les \( V_n^c\) sont fermés.
\end{proof}

\begin{corollary}       \label{CORooTWFYooCNMieM}
	Si \( X\) est un espace topologique métrisable, alors \( X\) accepte une base dénombrable de topologie autour de chaque point.
\end{corollary}

\begin{proof}
	Il s'agit seulement de remarquer que les singletons sont fermés et d'appliquer la proposition~\ref{PROPooXWBTooCvGLOj}.
\end{proof}

%-------------------------------------------------------
\subsection{Connexité}
%----------------------------------------------------


\begin{proposition}[\cite{MonCerveau}]	\label{PROPooUQNTooAegKwA}
	Tout espace métrique\footnote{Définition \ref{DefMVNVFsX}.} est complètement normal.
\end{proposition}

\begin{proof}
	Soient un espace métrique \( (X,d)\) ainsi que deux parties \( A,B\subset X\) telles que \( A\cap \bar B=B\cap\bar A=\emptyset\). Vu que \( A\cap \bar B=\emptyset\), pour chaque \( a\in A\), il existe une boule \( B(a,r_a)\) telle que \( B(a,r_a)\cap B=\emptyset\). Même histoire pour \( b\in B\).

	Nous posons \( U=\bigcup_{a\in A}B(a,r_a/2)\) et \( V=\bigcup_{b\in B}B(b,r_b/2)\). Les parties \( U\) et \( V\) sont ouvertes comme union d'ouverts. Nous avons évidemment \( A\subset U\) et \( B\subset V\).

	Pour voir que \( U\cap V=\emptyset\), supposons \( x\in U\cap V\). Il existe \( a\in A\) et \( b\in B\) tel que \( x\in B(a,r_a/2)\cap B(b,r_b/2)\). Remarquons que \( d(a,b)>r_a\) et \( d(a,b)>r_b\) : nous avons choisi \( r_a\) et \( r_b\) pour cela. Nous avons le calcul
	\begin{equation}
		d(a,b)  \leq d(a,x)+d(x,b)
		\leq \frac{ r_a+t_b }{2}
		<d(a,b).
	\end{equation}
	Cela est une contradiction.
\end{proof}

\begin{proposition}[\cite{MonCerveau}]	\label{PROPooBMFBooIQJMNC}
	Tout o-connexe\footnote{Définition \ref{DEFooJYCDooGNEDoC}} d'un espace métrique est connexe\footnote{Définition \ref{PROPooNLXDooVTZLan}}.
\end{proposition}

\begin{proof}
	Une espace métrique est complètement normal (proposition \ref{PROPooUQNTooAegKwA}) et tout o-connexe d'un espace complètement normal est connexe (proposition \ref{PROPooAHWNooXhrqnd}).
\end{proof}

%---------------------------------------------------------------------------------------------------------------------------
\subsection{Convexité}
%---------------------------------------------------------------------------------------------------------------------------

\begin{definition}[Partie convexe]        \label{DEFooQQEOooAFKbcQ}
	Une partie \( A\) d'un espace vectoriel est \defe{convexe}{partie convexe} si pour tout \( a,b\in A\) et pour tout \( t\in \mathopen[ 0 , 1 \mathclose]\), le point \( ta+(1-t)b\) est dans \( A\).

	Autrement dit, une partie est convexe lorsqu'elle contient tous les segments joignant ses points.
\end{definition}

\begin{proposition}[\cite{BIBooFHMTooBpSuHQ}] \label{PROPooJOCEooUKhkqQ}
	Toute intersection de convexes est convexe.
\end{proposition}

\begin{proof}
	Soit un espace vectoriel \( E\) ainsi que des parties convexes \( \{  C_i \}_{i\in I}\) indexées par un ensemble quelconque \( I\). Nous prouvons que \( C= \bigcap_{i\in I}C_i\) est convexe.

	Soient \( x, y\in C\), ainsi que \( i\in I\). Nous avons \( x,y\in C_i\) et donc \( \{ tx+(1-t)y \}_{t\in \mathopen[ 0 , 1 \mathclose]}\subset C_i\). Vu que cela est vrai pour tout \( i\), nous avons
	\begin{equation}
		\{ tx+(1-t)y \}_{t\in \mathopen[ 0 , 1 \mathclose]}\subset \bigcap_{i\in I}C_i,
	\end{equation}
	et donc le résultat attendu.
\end{proof}


\begin{proposition}[\cite{TQSWRiz}]	\label{PROPooJOXUooEZELna}
	Si \( A\) et \( B\) sont des convexes d'un espace vectoriel topologique, alors \( A-B\) est convexe.
\end{proposition}

\begin{proof}
	Soient \( a_1,a_2\in A\) et \( b_1,b_2\in B\). Nous devons prouver que \( t(a_1-b_1)+(1-t)(a_2-b_2)\) est dans \( A-B\). C'est un simple calcul :
	\begin{equation}
		t(a_1-b_1)+(1-t)(a_2-b_2)=\underbrace{ta_1+(1-t)a_2}_{\in A}-\underbrace{\big( tb_1+(1-t)b_2 \big)}_{\in B}\in A-B.
	\end{equation}
\end{proof}

\begin{lemma}		\label{LEMooEFCCooOuStrb}
	L'intérieur d'un convexe est convexe.
	%TODOooKTMWooQQLDkq. Prouver ça.
\end{lemma}

\ssdem

%---------------------------------------------------------------------------------------------------------------------------
\subsection{Norme}
%---------------------------------------------------------------------------------------------------------------------------

\begin{definition}[\cite{BrunelleMatricielle}, thème~\ref{THEMEooUJVXooZdlmHj}]  \label{DefNorme}
	Soit \( E\) un espace vectoriel (pas spécialement de dimension finie) sur le corps \( \eK\) (\( =\eR\) ou \( \eC\)). Une  \defe{norme}{norme} sur \( E\) est une application \( N\colon E\to\mathopen[ 0 , \infty \mathclose[\) telle que
	\begin{enumerate}
		\item
		      \( N(x)\geq 0\)
		\item
		      \( N(x)=0\) si et seulement si \( x=0\);
		\item   \label{ItemDefNormeii}
		      \( N(\lambda x)=| \lambda |N(x)\)
		\item   \label{ItemDefNormeiii}
		      \( N(x+y)\leq N(x)+N(y)\)
	\end{enumerate}
	pour tout \( x,y\in E\) et pour tout \( \lambda\in\eK\).

	La propriété~\ref{ItemDefNormeiii} est appelée \defe{inégalité triangulaire}{inégalité!triangulaire}.

	Un espace vectoriel muni d'une norme est un \defe{espace vectoriel normé}{espace vectoriel normé}.
\end{definition}
En prenant \( \lambda=-1\) dans la propriété~\ref{ItemDefNormeii}, nous trouvons immédiatement que \( N(-x)=N(x)\).

\begin{proposition}		\label{PropNmNNm}
	Toute norme \( N\) sur l'espace vectoriel \( E\) vérifie l'inégalité
	\begin{equation}
		\big| N(x)-N(y) \big|\leq N(x-y)
	\end{equation}
	pour tout \( x,y\in E\).
\end{proposition}

\begin{proof}
	Nous avons, en utilisant le point~\ref{ItemDefNormeiii} de la définition~\ref{DefNorme},
	\begin{subequations}
		\begin{align}
			N(x) & =N(x-y+y)\leq N(x-y)+N(y),	\label{subEqNNNxxyyya} \\
			N(y) & =N(y-x+x)\leq N(y-x)+N(x).	\label{subEqNNNxxyyyb}
		\end{align}
	\end{subequations}
	Supposons d'abord que \( N(x)\geq N(y)\). Dans ce cas, en utilisant \eqref{subEqNNNxxyyya},
	\begin{equation}
		\big| N(x)-N(y) \big|=N(x)-N(y)\leq N(x-y)+N(y)-N(y)=N(x-y).
	\end{equation}
	Si par contre \( N(x)\leq N(y)\), alors nous utilisons \eqref{subEqNNNxxyyyb} et nous trouvons
	\begin{equation}
		\big| N(x)-N(y) \big|=N(y)-N(x)\leq N(y-x)+N(x)-N(x)=N(y-x)=N(x-y).
	\end{equation}
	Dans les deux cas, nous avons retrouvé l'inégalité annoncée.
\end{proof}
Cette proposition signifie aussi que
\begin{equation}	\label{EqNleqNNleqNvqlqbs}
	-N(x-y)\leq N(x)-N(y)\leq N(x-y).
\end{equation}

Le lemme suivant dit que nous pouvons remplacer l'inégalité triangulaire par la convexité de la boule unité dans la définition de norme.
\begin{lemma}[\cite{BIBooXHRKooZrKDcs}]     \label{LEMooAVIJooFhdXXr}
	Soit une application \( N\colon E\to \mathopen[ 0 , \infty \mathclose[\) telle que
	\begin{enumerate}
		\item
		      \( N(x)\geq 0\) pour tout \( x\in E\),
		\item
		      \( N(x)=0\) si et seulement si \( x=0\),
		\item
		      \( N(\lambda x)=| \lambda | N(x)\).
	\end{enumerate}
	Alors \( N\) est une norme si et seulement si la partie
	\begin{equation}
		B = \{ x\in E\tq N(x)\leq 1 \}
	\end{equation}
	est convexe\footnote{Définition \ref{DEFooQQEOooAFKbcQ}.}.
\end{lemma}

\begin{proof}
	Dans les deux sens.
	\begin{subproof}
		\spitem[\( \Rightarrow\)]
		Nous supposons que \( N\) est une norme et nous prouvons que la boule \( B\) est convexe. Soient \( x,y\in B\) et \( \lambda\in \mathopen[ 0 , 1 \mathclose]\). Nous avons
		\begin{subequations}
			\begin{align}
				N\big( \lambda x+(1-\lambda)y \big) & \leq N(\lambda x)+N\big( (1-\lambda)y \big) \\
				                                    & = \lambda N(x)+(1-\lambda)N(y)              \\
				                                    & \leq \lambda +(1-\lambda)                   \\
				                                    & = 1.
			\end{align}
		\end{subequations}
		Nous avons utilisé diverses propriétés de la norme, ainsi que la majoration \( N(x), N(y)\leq 1\).
		\spitem[\( \Leftarrow\)]
		Nous supposons que \( B\) est convexe, et nous prouvons que \( N\) vérifie l'inégalité triangulaire. Soient \( x,y\in E\) que nous choisissons tous deux non nuls (sinon c'est trop facile). Nous posons
		\begin{equation}        \label{EQooCIMBooFeOtWg}
			z=\frac{ x+y }{ N(x)+N(y) }
		\end{equation}
		et la subtilité sera d'écrire \( z\) de telle sorte à être une somme de deux éléments de \( B\). L'astuce est de poser
		\begin{equation}
			\lambda=\frac{ N(x) }{ N(x)+N(y) }.
		\end{equation}
		Une simple vérification montre qu'alors
		\begin{equation}
			z = \lambda\frac{ x }{ N(x) }+  (1-\lambda) \frac{ y }{ N(y) }.
		\end{equation}
		Nous avons évidemment \( x/N(x)\in B\) (et de même avec \( y\)). Puisque \( B\) est convexe, nous avons \( z\in B\). Exprimons le fait que \( z\in B\) à partir de la définition \eqref{EQooCIMBooFeOtWg} :
		\begin{equation}
			\frac{ N(x+y) }{ N(x)+N(y) }\leq 1.
		\end{equation}
		Cela signifie exactement \( N(x+y)\leq N(x)+N(y)\).
	\end{subproof}
\end{proof}

\begin{normaltext}
	Afin de suivre une notation proche de celle de la valeur absolue, à partir de maintenant, la norme d'un vecteur \( v\) sera notée \( \| v\|\) au lieu de \( N(v)\). La proposition~\ref{PropNmNNm} s'énoncera donc
	\begin{equation}
		\big| \| x \|-\| y \| \big|\leq \| x-y \|.
	\end{equation}
	Un espace vectoriel \( E\) muni d'une norme est, on l'a déjà dit, un \defe{espace vectoriel normé}{normé!espace vectoriel}; on le notera \( (E,\| . \|)\) pour distinguer la norme fixée.
\end{normaltext}

Une autre inégalité utile de temps en temps.
\begin{corollary}       \label{CORooDFBGooAqVRfS}
	Si \( a\) et \( b\) sont dans un espace vectoriel normé, alors
	\begin{equation}
		\big| \| a-b \|-\| b \| \big|\leq \| a \|.
	\end{equation}
\end{corollary}

\begin{proof}
	Il s'agit seulement de la proposition \ref{PropNmNNm} avec \( x=a-b\) et \( y=b\).
\end{proof}

\begin{lemmaDef}[Distance induite par une norme]        \label{LEMooWGBJooYTDYIK}
	Soit un espace vectoriel normé \( (E,\| . \|)\). Nous posons
	\begin{equation}        \label{EQooZYJRooAHnsIG}
		d(x,y)=\| x-y \| .
	\end{equation}
	Alors
	\begin{enumerate}
		\item       \label{ITEMooLITDooPeReOk}
		      \( d\) est invariante par translations : \( d(a,b)=d(a+u,b+u)\)
		\item
		      \( d\) est une distance\footnote{Définition~\ref{DefMVNVFsX}.} sur \( E\).
	\end{enumerate}
	C'est la \defe{distance induite}{distance!associée à une norme} par la norme.
\end{lemmaDef}

\begin{proof}
	Le fait que la formule \eqref{EQooZYJRooAHnsIG} soit invariante par translations est immédiat. En ce qui concerne le fait que ce soit une distance, le seul point délicat à vérifier est l'inégalité triangulaire. Mais, pour tous \( x, y, z \in E\), on a
	\begin{equation}
		d(x,y)=\| x-y \| = \| x-z+z-y \|  \leq\| x - z \|+\| z - y\| =d(x,z)+d(z,y).
	\end{equation}
\end{proof}

%-------------------------------------------------------
\subsection{Topologie}
%----------------------------------------------------

\begin{definition}      \label{DEFooPMVFooPSYVNQ}
	La topologie associée à une norme est celle associée à la distance donnée en \ref{LEMooWGBJooYTDYIK} par le théorème \ref{ThoORdLYUu}.
\end{definition}
